\subsection{Which quotients retain the Selmer argument}

All derived and lower central subgroups in this part are closed. A section
of a geometric quotient is called Selmer only when its localization is the
\emph{image} of a local point or cusp section. This condition is stronger
than satisfying a selection of finite local obstruction equations.

\begin{statement}{B.1}{Finite support for a maximal prosolvable quotient}{Proved}
Let $U/\mathbb Q$ be a smooth geometrically connected affine curve, let
$X$ be its smooth compactification, and suppose $g(X)\geq1$. Put
$P=\pi_1(U_{\overline{\mathbb Q}})$ and
$P^{(\infty)}=\bigcap_{j\geq1}P^{(j)}$. A Selmer section of the
geometrically maximal prosolvable arithmetic quotient whose underlying
local points belong to a finite reduced subscheme $Z\subset X$ at a
density-one set of finite primes belongs to the image of a point or cusp
decomposition group above $Z(\mathbb Q)$.

In particular, the same conclusion holds if $X$ has a nonconstant
morphism to a proper curve $W$ satisfying
$W(\mathbb A_{\mathbb Q})_\bullet^{\mathrm{f\text{-}ab}}=W(\mathbb Q)$.
\end{statement}
\begin{proof}
The essential point is closure of the neighbourhood category under
finite abelian refinements. Let $H\subset P$ be open with
$P^{(\infty)}\subset H$. The finite quotient by the core of $H$ is
solvable, so $P^{(m)}\subset H$ for some $m$. Consequently
\[
 P^{(\infty)}\subset P^{(m+1)}\subset H'.
\]
Every finite abelian geometric cover of this neighbourhood therefore
still belongs to the same prosolvable category. Local point and cusp
sections lift along a neighbourhood by the finite-fibre argument of A.1.
After compactification the resulting local points survive all finite
abelian descent. This is precisely the abelian variant of
\citet[Theorem 2.1 and its subsequent variant for characteristic geometric
quotients]{AMHarariStix2012}; here the geometric quotient used is the one
just checked to contain every finite abelian refinement.

For a neighbourhood $U'$, let $X'$ be its compactification and
$Z'=(X'\times_X Z)_{\mathrm{red}}$. Riemann--Hurwitz gives $g(X')\geq1$.
The hypotheses of \citet[Theorem 8.2]{AMStoll2007} hold for $X',Z'$:
the adelic point survives finite abelian descent and has support in
$Z'$ on a density-one set. Hence $Z'(\mathbb Q)\ne\varnothing$.
The sets $Z'(\mathbb Q)$ are finite. A common refinement supplies
compatible points for any finite collection of neighbourhoods, so their
inverse limit is nonempty. The corresponding pro-point is fixed by the
section. Its stabilizer in the prosolvable arithmetic quotient is the
image of the appropriate point or cusp decomposition group. This proves
the first assertion.

In the second assertion functoriality of finite abelian descent gives a
single $w\in W(\mathbb Q)$ as the image of all finite local points.
The fibre $Z=f^{-1}(w)_{\mathrm{red}}$ is finite, since $f$ is a
nonconstant morphism of proper curves, and the first assertion applies.
The modified adelic space is used at the real place throughout; no
choice of a real point inside a connected component is asserted.

The argument is valid over any number field. It also permits the
target to be found after a finite extension $F/\mathbb Q$.
Indeed, suppose $X_F$ has a nonconstant map to a finite-abelian
excellent proper curve. Applying the argument over $F$ gives a
point $x\in X(F)$ supporting the restricted section. At a finite
place $w$ of $F$ above $p$, its compact abelianization equals that
of the original local point $P_p$, viewed over $F_w$.
The full Tate-module Kummer map on the compact group
$J_X(F_w)$ is injective. The Abel--Jacobi map
$y\mapsto[y-x]$ is injective for a positive-genus curve, so
$P_p=x$ over $F_w$. Thus every $P_p$ lies in the finite reduced
$\mathbb Q$-scheme consisting of the conjugates of $x$.
Applying the first assertion over $\mathbb Q$ gives a rational
point or cusp section. In particular, if $X(\mathbb Q)$ is
contained in the boundary, the original section is cuspidal.
This descent uses no assertion about normalizers of cusp
decomposition groups in the prosolvable quotient.
\end{proof}

\begin{statement}{B.2}{The quotient in Lead 1}{Proved}
Let $\ell\geq5$, let $P=\pi_1(U_{\ell,\overline{\mathbb Q}})$,
let $K$ be the kernel of the mod-$\ell$ Legendre representation on $P$,
and put $R=K^{(\infty)}$. Every Selmer section of the arithmetic
quotient with geometric group $P/R$ belongs to the image of a rational
cusp decomposition group.
\end{statement}
\begin{proof}
The residual representation is defined on this quotient because
$R\subset K$. The corrected local and modular argument of Part A for
R3 uses only the residual representation and the local point or cusp
descriptions. It therefore applies verbatim to this quotient section:
the residual representation preserves a line. Choose that line and
let $H$ be the geometric subgroup of the resulting connected component
of $U_\ell\times_{\mathbb P^1-\{0,1,\infty\}} Z_\ell$. Then
$K\subset H$ and $H/K$ is a subgroup of a Borel group and is solvable.
If its derived length is $d$, $H^{(d)}\subset K$. The two inclusions
$K\subset H$ and $H^{(d)}\subset K$ give
\[
 K^{(\infty)}\subset H^{(\infty)}
 \quad\hbox{and}\quad
 H^{(\infty)}\subset K^{(\infty)};
 \qquad H^{(\infty)}=R.
\]
Thus the lifted section is a Selmer section of the maximal prosolvable
quotient of this neighbourhood. Its compactification maps
nonconstantly to $X_0(4\ell)$, of positive genus. The precise strength
of \citet[Corollary 8.8]{AMStoll2007} is finite \emph{abelian}
excellence, which is sufficient for B.1. We obtain a rational point or
cusp upstairs. A point in the open curve would give, after the quadratic
twist and possibly the $\ell$-isogeny in Part A, an elliptic curve over
$\mathbb Q$ with full rational two-torsion and a rational point of
order $\ell$. Mazur's rational-torsion theorem excludes it. The point
upstairs is consequently a cusp, and its image is a rational cusp of
$U_\ell$. No lifting from $P/R$ to $P$ has been assumed.
\end{proof}

\begin{statement}{B.3}{The cubic and classical quartic prosolvable quotients}{Proved}
Every Selmer section of the geometrically maximal prosolvable quotient
of $U_3$ or of the classical Fermat quartic $C_4$ is cuspidal.
\end{statement}
\begin{proof}
The compact cubic is $X_0(27)$. The classical quartic has the
nonconstant map to the rank-zero CM elliptic curve used and verified in
Part A. The targets there are finite-abelian excellent: for $X_0(27)$
this is \citet[Corollary 8.8]{AMStoll2007}; for the quartic target it
is \citet[Theorem 8.6]{AMStoll2007}, with finite Mordell--Weil group
and finite Tate--Shafarevich group checked in Part A. Both source
compactifications have positive genus. B.1 therefore gives a rational
point or cusp. The rational-point calculations in Part A place every
rational point on the boundary. This proves the assertion for the
\emph{classical} quartic, without substituting the real-point-free plus
quartic.
\end{proof}

\begin{statement}{B.4}{The septic maximal prosolvable quotient}{Proved}
Every Selmer section of the geometrically maximal prosolvable quotient
of $U_7$ is cuspidal. An explicit target for B.1 is $X_0(49)$.
\end{statement}
\begin{proof}
Write $a^7+b^7+c^7=0$. The formulas
\[
 (X,Y,Z)=(a^3c,b^3a,c^3b)
\]
give a map to the Klein quartic
$X^3Y+Y^3Z+Z^3X=0$, because the left side pulls back to
$a^3b^3c^3(a^7+b^7+c^7)$. With
$e_1=X+Y+Z$, $e_2=XY+YZ+ZX$ and $e_3=XYZ$, put
\[
 x=-\frac{e_2^2}{e_1e_3},\qquad
 y=\frac{2e_1^2e_2-5e_2^2}{e_1e_3}.
\]
These functions satisfy $y^2=4x^3+21x^2+28x$. Indeed the numerator
of the difference, after multiplication by $(e_1e_3)^3$, is
\[
 4e_2^2(X^3Y+Y^3Z+Z^3X)(X^3Z+Z^3Y+Y^3X).
\]
Both maps are nonconstant: the first has finite generic fibres, and
the second is the quotient by the cyclic coordinate permutation of
order three. A rational map between smooth proper curves extends over
the omitted finitely many points. Setting $u=x+2$ and $v=(y-u)/2$
gives the minimal equation
\[
 E_{49}:\quad v^2+uv=u^3-u^2-2u-1.
\]
The identification $E_{49}\simeq X_0(49)$, including this equation
and the maps from the Klein quartic, is established in
\citet[\S\S2.1, 3.2, 4.2]{BElkies1999}. Stoll's Corollary 8.8
therefore supplies finite-abelian excellence. Apply B.1 and the
rational-point conclusion for the Fermat septic in Part A.

The program \texttt{b\_septic.py} checks the displayed polynomial
identities exactly. It also gives a table-independent nonvanishing
certificate for $L(E_{49},1)$. For completeness its analytic input is
as follows. The unique normalized weight-two form at level $49$ has
coefficients obtained from the minimal equation by point counting and
the Hecke recurrences. Since the Fricke involution fixes $i/7$ and is
nonidentity, it acts by $-1$ on the differential of this genus-one
curve. The sign of the functional equation is $+1$, and splitting the
Mellin integral at $1/7$ gives
\[
 L(E_{49},1)=2\sum_{m\geq1}\frac{a_m}{m}e^{-2\pi m/7}.
\]
The bound $|a_m|\leq d(m)\sqrt m\leq2m$ bounds the omitted tail
after $m=30$ by $4r^{31}/(1-r)$, $r=e^{-2\pi/7}$. Rational bounds
$333/106<\pi<355/113$ and alternating Taylor bounds for $e^{-t}$
give the outward displayed interval
$[0.966654299,0.966683148]$; all actual comparisons in the program
are rational, and the tail is less than $5.563\cdot10^{-12}$.
The extra rank-zero check uses the $2$-isogeny models
$Y^2=X^3+21X^2+112X$ and $Y^2=X^3-42X^2-7X$.
The excluded covering squareclasses $2$ and $-1$ have, respectively,
primitive quartic residues $\{2,6,7,8,12,15\}$ and
$\{7,12,15\}$ modulo $16$, disjoint from the squares.
The two descent images have orders two and two, so the rank is zero.
The argument using Stoll's modular corollary does not depend on this
additional computation.
\end{proof}

\begin{statement}{B.5}{The least positive derived length for the cubic is one}{Proved}
Every Selmer section of the \emph{full profinite geometric
abelianization} of $U_3$ is cuspidal. Thus the least positive integer
$k$ asked for in B.3 of the prompt is $k=1$ for $n=3$.
\end{statement}
\begin{proof}
Fix the cusp $O$ as tangent origin. Let $E=X_3$, let $D=E[3]$, and
write $V=TJ_{E,D}$ for the full profinite Tate module of the generalized
Jacobian. There is an exact sequence
\[
 0\longrightarrow W\longrightarrow V\longrightarrow TE
 \longrightarrow0,\qquad
 W=\widehat{\mathbb Z}(1)[D]/\widehat{\mathbb Z}(1)\sum_{d\in D}d.
\]
The compact projection of a Selmer class is a full Tate-module Selmer
class for $E$. The rank-zero and finite-$\Sha$ facts proved in Part A
identify this group with $E(\mathbb Q)$, consisting of the three
rational cusps. For any finite place $v$, the full local Kummer map
$E(\mathbb Q_v)\to H^1(\mathbb Q_v,TE)$ is injective: a compact
$p$-adic Lie group has trivial intersection of all its multiplication
subgroups. Hence every local underlying point is the same rational
cusp. Relabel it $O$ and subtract its chosen tangent section.

The resulting class $z\in H^1(\mathbb Q,V)$ has zero compact
projection, hence comes from $H^1(\mathbb Q,W)$. There is no ambiguity
from $H^0(\mathbb Q,TE)$, which is zero because the rational torsion
group is finite. Locally $z$ belongs to the line $I_O$ generated by
the cusp inertia. Thus its image in $H^1(\mathbb Q,W/I_O)$ is locally
zero. The boundary has three rational points and three conjugate
pairs over $K=\mathbb Q(\zeta_3)$. Eliminating the $O$ coordinate and
using a second rational cusp to split the remaining diagonal gives
the integral permutation-module isomorphism
\[
 W/I_O\simeq\widehat{\mathbb Z}(1)
 \oplus\bigl(\operatorname{Ind}_{G_K}^{G_{\mathbb Q}}
                  \widehat{\mathbb Z}(1)\bigr)^3.
\]
For $F=\mathbb Q$ and $F=K$, the map
$H^1(F,\widehat{\mathbb Z}(1))\to\prod_w
H^1(F_w,\widehat{\mathbb Z}(1))$ is injective. To see this directly,
both fields have class number one and finite unit group. Their
completed multiplicative groups have prime-valuation coordinates and
a finite root-of-unity coordinate. Local valuations detect all the
prime coordinates; a remaining root of unity is detected at a local
place of matching residue characteristic. Shapiro's lemma now proves
the required injectivity for $W/I_O$. It follows that $z$ comes from
$H^1(\mathbb Q,I_O)$, exactly the cusp packet.

This proves the least-positive-$k$ assertion literally as asked.
It gives no monotonicity assertion in $k$: a higher quotient can have
additional classes in the fibre over a cuspidal abelian class. In
particular, finding noncuspidal classes at $k=2$ does not contradict
the answer $k=1$.
\end{proof}

\begin{statement}{B.6}{The zero-abelian metabelian families are not Selmer}{Proved}
In the full profinite metabelian families constructed in LM6 E.31,
E.33 and E.34, every nonzero parameter is non-Selmer; the zero
parameter is the specified cusp section. The obstruction to being
Selmer occurs at the coefficient prime, respectively $3$ and $5$.
\end{statement}
\begin{proof}
We state explicitly the features of the family needed here. Its
abelianization is the base tangent class at a rational cusp $O$;
its first nonzero central class lies in
$H^1(\Gamma,\Lambda^2 V_\ell)$, and the parameter map into that
group is injective. Here $\Gamma=P_{\mathrm{ar}}\rtimes C_{\ell-1}$,
where $P_{\mathrm{ar}}$ is the maximal pro-$\ell$ Galois group over
$K=\mathbb Q(\zeta_\ell)$ unramified away from $\ell$. These are
exactly the detecting classes in the definitions of the cited
families, rather than an assertion about every possible metabelian
section.

The decomposition group at the unique prime above $\ell$ surjects
onto $P_{\mathrm{ar}}$. If it did not, its image would be contained
in a maximal open subgroup, giving a cyclic degree-$\ell$ extension
of $K$, unramified away from $\ell$, in which that prime splits
completely. This would be everywhere unramified. The fields
$\mathbb Q(\zeta_3)$ and $\mathbb Q(\zeta_5)$ have class number one,
contradicting class field theory. The local cyclotomic extension also
contains the entire $C_{\ell-1}$. Thus $G_{\mathbb Q_\ell}$ surjects
onto $\Gamma$, and inflation gives an injection
\[
 H^1(\Gamma,\Lambda^2 V_\ell)
 \lhook\joinrel\longrightarrow
 H^1(\mathbb Q_\ell,\Lambda^2 V_\ell).
\]
One can check injectivity without any finiteness hypothesis: a
local coboundary defined by a vector is a coboundary on $\Gamma$
because the map to $\Gamma$ is surjective.

If a family member were locally geometric at $\ell$, its full
profinite compact abelianization and injectivity of the local full
Tate-module Kummer map would force its underlying point to be $O$.
We use here the Abel--Jacobi embedding of the positive-genus
compactification; equality of its divisor class with $O$ implies
equality of the points. The local cusp parameter must then be zero.
Indeed $H^0(\mathbb Q_\ell,V/I_O)=0$: after filling the compact
part its possible invariants lie in local Tate modules of roots of
unity, all zero, and the compact Tate invariants are zero since local
torsion is finite. Thus
$H^1(\mathbb Q_\ell,I_O)\to H^1(\mathbb Q_\ell,V)$ is injective.
The localized section would be the base cusp class itself, so its
central detecting class would vanish. The displayed injection and
the injectivity of the family parameter then force the parameter to
be zero. This proves the claim, including all $3^r$ and $5^r$ in
these specified full profinite metabelian families.
\end{proof}

\subsection{An infinite-ramification counterexample with a continuous lift}

\begin{statement}{B.7}{The class-three form of B.11 is false}{Refuted}
There is an explicit noncuspidal class in
$H^1(\mathbb Q,N/\gamma_3N)$, for the rank-ten free pro-$3$
geometric group of the nine-punctured cubic, ramified at infinitely
many primes, which lifts to $H^1(\mathbb Q,N/\gamma_4N)$.
It has zero abelianization. The construction does not require the
unevaluated Ross coordinate of the arithmetic action.
\end{statement}
\begin{proof}
We first specify an integral Tate vector independently of the inherited
action matrices. Identify $U_3=E\setminus E[3]$. The isogeny $[3]$
gives a degree-nine cover
$E\setminus E[3]\to E\setminus\{O\}$. Geometric transfer gives a
$G_{\mathbb Q}$-equivariant map
\[
 j:T_3E=\pi_1(E_{\overline{\mathbb Q}}\setminus\{O\})^{(3),ab}
       \longrightarrow V=N^{ab}.
\]
Transfer followed by the cover pushforward is multiplication by nine;
the cover pushforward on compact Tate modules is multiplication by
three. Therefore the compact projection of $j$ is $3\operatorname{id}$.
Choose a symplectic basis $e,f$ of $T_3E$ and put
$w=j(e)\wedge j(f)\in\Lambda^2V$. Then $w$ is integral, nonzero,
and transforms by $\chi_3$. Its compact projection is nine times the
determinant generator.

For full finite coordinates, take the Schreier graph of the kernel
$F_2\to(\mathbb Z/3)^2$ with vertex representatives $x^iy^j$,
$0\leq i,j<3$, and tree edges $(i,0,x)$ for $i<2$ and $(i,j,y)$
for $j<2$. Order the ten remaining edges as
\[
 (01x,02x,02y,11x,12x,12y,20x,21x,22x,22y).
\]
In this basis the two transfer vectors are
\[
 (1,1,0,1,1,0,1,1,1,0),\qquad
 (0,0,1,0,0,1,0,0,0,1),
\]
and $w$ is their exterior product. This specifies all 45 coordinates,
without a choice of the action cocycle.

Set $A=\gamma_3N/\gamma_4N=\operatorname{Lie}_3(V)$ and
$L=\gamma_2N/\gamma_4N$. The module $L$ is abelian, since
$[\gamma_2N,\gamma_2N]\subset\gamma_4N$, and fits in
\[
 0\longrightarrow A\longrightarrow L\longrightarrow
 \Lambda^2V\longrightarrow0.                                      \tag{B.7.1}
\]
Choose a lift $\omega$ of $w$, made anti-invariant under complex
conjugation by averaging with $1/2$, and define
\[
 \eta(g)=\chi_3(g)^{-1}\bigl(g\omega-\chi_3(g)\omega\bigr)
       \in A(-1).
\]
This is an integral continuous cocycle, unramified outside $3$:
the punctured curve with its unit tangent has good reduction away
from $3$. For a Kummer class $\beta\in H^1(\mathbb Q,\mathbb Z_3(1))$,
the zero-abelian class $\beta w$ has obstruction $[\eta\cup\beta]$
to lifting along (B.7.1). This follows directly by taking
$g\mapsto\beta(g)\omega$ as a cochain lift; its defect is
$\chi_3(g)\beta(h)\eta(g)$.

Enumerate the rational primes $p_j\equiv2\pmod3$ increasingly, starting
with $2$, and put
\[
 h_j=v_3(p_j+1),\qquad
 a_j=3^{j+330h_j+128},\qquad
 \beta=\sum_{j\geq1}a_j\kappa_3(p_j).                              \tag{B.7.2}
\]
The sum converges uniformly as a sum of integral continuous cocycles.
At $p_j$ its inertia value is $a_j$, and it is zero on inertia at
every other prime different from $3$.

Here is the exact global lifting check. Work over
$K=\mathbb Q(\zeta_3)$, with $S$ the prime above $3$. At the prime
$\mathfrak p$ above inert $p$ the residue cardinality is $q=p^2$.
The eigenvalues of Frobenius on $V_{\mathbb Q_3}$ are $-p$ twice
and $p^2$ eight times: the compact elliptic factor has $a_p=0$
(cubing is bijective on $\mathbb F_p$, so the cubic has $p+1$
points), and the eight boundary coordinates are Tate coordinates over $K$.
The multigraded Witt dimensions for $A$ are
$2,32,128,168$. Consequently the eigenvalues of
$q^{-1}\Phi_A$ are
\[
 -p\ (2),\qquad p^2\ (32),\qquad -p^3\ (128),\qquad p^4\ (168).
\]
For $D_p=q^{-1}\Phi_A-1$, the lifting-the-exponent formulas give
\[
 v_3(\det D_p)=2h+32h+128(h+1)+168h=330h+128.                     \tag{B.7.3}
\]
The operator $D_p$ is an integral matrix, invertible over $\mathbb Q_3$.
Its adjugate is integral. Therefore the unique rational tame correction
\[
 c_{\mathfrak p_j}=-a_jD_{p_j}^{-1}\eta(F_{\mathfrak p_j})
\]
belongs to $3^jA(-1)$. All other corrections are zero. The
corrections are thus integral and tend to zero, which is stronger than
merely vanishing at each finite coefficient level.

We justify the sufficiency of this check rather than importing it
from LM7 E.44. For $G_{K,S}$, Kummer theory and the class number one
of $K$ give $\dim_{\mathbb F_3}H^1(G_{K,S},\mathbb F_3)=2$.
The global Euler characteristic
\citep[I, Theorem 5.1]{BMilne2006}, with one complex place, gives
$H^2(G_{K,S},\mathbb F_3)=0$. The mod-$3$ action on $V$ is
unipotent: $E[3]$ and all boundary points are rational over $K$.
Hence every composition factor of $A/3^mA$ is trivial, and induction
in a composition series gives $H^2(G_{K,S},A/3^mA)=0$ for all $m$.
These finite-support $H^1$ groups are finite, so their inverse system
is Mittag--Leffler and
$H^2_{\mathrm{cts}}(G_{K,S},A)=0$. Also $H^3$ vanishes; there are
no real places and $3$ is invertible on the corresponding open
arithmetic curve. The identification of Galois with etale cohomology
is \citet[II, Proposition 2.9]{BMilne2006}.

Apply the localization triangle with finite coefficients and then its
derived inverse limit, as in \citet[\S3.6]{BJannsen1988}. The residue
complex is the inverse limit of the direct sums of two-term complexes
\[
 \bigoplus_{\mathfrak q\notin S}
 [\,A/3^mA(-1)\xrightarrow{\Phi_{\mathfrak q}/q-1}
                       A/3^mA(-1)\,].
\]
Its inverse-limit terms are the modules of families tending to zero.
The finite-support cohomology just proved to vanish implies that a
global obstruction is zero precisely when its residue vector is the
image of such a restricted integral family. The family $c$ above
has exactly this property. Restriction from $\mathbb Q$ to $K$ is
injective in degree two for $3$-primary coefficients, by
restriction--corestriction and $[K:\mathbb Q]=2$. Thus the obstruction
over $\mathbb Q$ also vanishes, and (B.7.2) lifts to class three.

It remains to distinguish this from a cusp. Its inertia at every
$p_j$ is the nonzero central element $a_jw$, so it is ramified at
infinitely many primes. Its abelian class is zero. The other two
rational cusps have nonzero compact $3$-adic Kummer classes, since
their differences from $O$ are nonzero points of order three.
At $O$ the cusp-inertia map into $V$ induces an injection on global
$H^1$: $H^0(\mathbb Q,V/I_O)=0$ by good Frobenius weights.
Consequently the only cuspidal class with zero abelianization is the
base tangent itself, whose class-two inertia is zero. Our class is
not cuspidal.

The script \texttt{b\_infinite\_class3.py} constructs the stated
Schreier transfer vectors, verifies the multigraded dimensions, and
checks (B.7.3) at all 87 inert primes below 1000. The proof for all
primes is the displayed eigenvalue and valuation calculation, not an
extrapolation from that finite check.
\end{proof}

\subsection{The compact quartic action and the cyclotomic descent}

\begin{statement}{B.8}{The rank-two CM carrier and the factor two}{Known}
For the plus Fermat quartic, the pro-$5$ compact class-two outer
action class is carried by the rank-two CM module asserted in LM7
E.109. With the ordinary Magnus action convention, \emph{twice} the
action class corresponds to the Abel--Jacobi class of the Ceresa cycle.
The equality with Ceresa is the existing Hain--Matsumoto and
Bisogno--Li--Litt--Srinivasan comparison; the rank-two specialization
and an integral coefficient projector follow from the Fermat
automorphisms as below. This statement does not certify the additional
nonvanishing calculation in LM7 E.109.
\end{statement}
\begin{proof}
Put $R=\mathbb Z_5$, $H=T_5J_C$, $Z=\Lambda^2H/Rq$,
$L=(\Lambda^3H)(-1)$, and
$\mathcal D=\operatorname{Hom}(H,Z)/\operatorname{ad}H$.
The ranks are $6,14,20,78$. Since $2$ is invertible, the actual
class-two surface group is the BCH group $H\oplus Z$ with law
$(x,z)(y,t)=(x+y,z+t+x\wedge y/2)$. Thus an arithmetic outer action
has a central off-diagonal class $[c]\in H^1(\mathbb Q,\mathcal D)$.
The symplectic Johnson map is
\[
 f(a\wedge b\wedge d)(x)=
 \langle a,x\rangle b\wedge d+
 \langle b,x\rangle d\wedge a+
 \langle d,x\rangle a\wedge b\pmod{Rq}.
\]
It satisfies $f(q\wedge h)=-\operatorname{ad}_h$, and induces
$\bar f:L/(q\wedge H)\hookrightarrow\mathcal D$.

Over $F=\mathbb Q(\zeta_8)$ the deck group
$\mathcal A=\mu_4^2$ acts on $H$. Its six characters, up to the
duality convention, are
\[
 (1,1),(1,2),(2,1),\quad -(1,1),-(1,2),-(2,1)
       \quad\hbox{in }(\mathbb Z/4)^2.
\]
Among triples of these six characters, precisely the all-first-three
and the all-last-three triples sum to zero. There are no invariant
vectors in $q\wedge H$. It follows that
\[
 (L/(q\wedge H))^{\mathcal A}
   = M=(\Lambda^3H^+\oplus\Lambda^3H^-)(-1),
       \qquad \operatorname{rank}_R M=2.
\]
Here $H^\pm$ are the two rank-three eigenspaces for the common CM
operator $i$, equivalently the two indicated character triples.
The projectors are integral because $16$ is a $5$-adic unit.
Galois either preserves these two sets of characters or interchanges
them, so $M$ and its projector are defined over $\mathbb Q$.

For a symplectic basis $(e_r,f_r)_{r=1}^3$ and a lift
$\widetilde D:H\to\Lambda^2H$, define
\[
 \Pi(D)=\frac13\pi_M\sum_{r=1}^3
       \bigl(e_r\wedge\widetilde D(f_r)
                         -f_r\wedge\widetilde D(e_r)\bigr).
\]
Changing the lift adds $q\wedge H$, which $\pi_M$ kills.
Tensor duality proves basis independence. Direct expansion gives
$\operatorname{Alt}f(\xi)=3\xi$ modulo $q\wedge H$ and
$\operatorname{Alt}(\operatorname{ad}_h)=-2q\wedge h$.
Thus $\Pi$ is a well-defined equivariant map
$\mathcal D\to M$ with $\Pi\bar f|_M=1$; all denominators are
units in $R$.

\citet[Theorem 3 and \S8]{BHainMatsumoto2005} and
\citet[Corollary 2.17 and Proposition 2.23]{BBisognoEtAl2023}
identify $2[c]$ with $\bar f_*\nu(C)$, where $\nu(C)$ is the
basepoint-independent Ceresa class in the primitive exterior cube.
Their Magnus map is $g\mapsto(x\mapsto g(\widetilde x)
\widetilde x^{-1})$. In the group ring the map
$u\mapsto(u-1)-(u-1)^2/2\pmod{I^3}$ gives exactly the BCH
coordinates above, so this comparison has the same sign and factor
two. Curve automorphisms preserve the outer action class, and hence
the Ceresa class. Averaging over $\mathcal A$ over $F$ therefore
places it in $M$. Restriction from $\mathbb Q$ to $F$ is injective
on these cohomology groups since the degree is four, a $5$-adic
unit. Consequently
\[
 \beta=2\Pi_*[c]\in H^1(\mathbb Q,M),\qquad
 2[c]=\bar f_*\beta,
\]
and $\beta|_{G_F}$ is the pure-CM projection of the Ceresa
Abel--Jacobi class. Filling punctures respects this assertion: in
any affine BCH coordinates the filling map has a linear correction
$J$, and its off-diagonal classes satisfy
$c^C=p_Bc^Aj+Jj-g(Jj)$. Changes of arithmetic section add an inner
term, killed by $\Pi$. This verifies the actual action comparison,
not just the associated graded representation.
\end{proof}

\begin{statement}{B.9}{Lead 9: cohomological dimension of the cyclotomic field}{Proved}
For $K_\infty=\mathbb Q(\mu_{3^\infty})$,
$\operatorname{cd}_3(G_{K_\infty})\leq1$.
Every finite embedding problem over this group with a $3$-group
kernel has a weak solution. More precisely, for the cubic
geometric pro-$3$ group $N$ and every closed
$G_{K_\infty}$-stable normal subgroup $J$, every continuous
section of $(N/J)\rtimes G_{K_\infty}$ lifts to a continuous
section of $N\rtimes G_{K_\infty}$, with its prescribed
quotient section unchanged. For every finite free continuous
$\mathbb Z_3[G_{K_\infty}]$-module $M$,
\[
                  H^2_{\mathrm{cts}}(G_{K_\infty},M)=0.
\]
In particular there is no finite-coefficient phantom in
degree two for these lattices over $K_\infty$.
\end{statement}
\begin{proof}
We prove the Brauer criterion, including its quantifier over algebraic
extensions. Let $L/K_\infty$ be algebraic and let
$b\in\operatorname{Br}(L)[3^m]$. A central simple algebra representing
$b$ descends to a finite number field $L_0\subset L$. Its finitely
many structure constants and a splitting of $3^m b$ descend
after enlarging $L_0$, so its descended class may and will be
chosen $3^m$-torsion. Its finitely many nonzero local invariants
then have $3$-power denominators.
For every finite place $v$ of $L_0$, the local degrees in
$L_0(\mu_{3^r})/L_0$ have unbounded $3$-parts. At $3$ this follows
from the totally ramified cyclotomic tower. Away from $3$ the
Frobenius image in $1+3\mathbb Z_3$ has infinite order: a positive
integer norm greater than one cannot be a $3$-adic root of unity.
Passing to the fixed finite extension $L_0$ changes degrees only by
a bounded factor. Choose $r$ so that all the nonzero local invariants
are killed by multiplication by the corresponding local degrees.
There are no real places in the resulting cyclotomic compositum.
The Albert--Brauer--Hasse--Noether sequence
\citep[VIII, Theorem 4.2]{BMilneCFT2020} implies that the class is
zero over $L_0(\mu_{3^r})$, and hence over $L$.
Thus $\operatorname{Br}(L)\{3\}=0$ for every algebraic extension
$L/K_\infty$. The equivalence in
\citet[Theorem 6.1.8]{BGilleSzamuely2006} proves the assertion about
cohomological dimension.

For clarity the lifting consequence uses genuine compatibility, not
an assertion that unrelated finite solutions compactify. A finite
$3$-group kernel has a nonzero characteristic elementary abelian
subgroup in its centre. Induction on the kernel order reduces a
finite embedding problem to abelian lifting problems, whose
obstructions are $H^2$ classes and vanish by the dimension bound.
Choose a nested cofinal sequence of characteristic finite quotients
by open subgroups $C_m$ of the finitely generated group $N$,
with $\bigcap_mC_m=1$ and $C_0=N$. The actual equivariant
maps
\[
 N/(J\cap C_{m+1})\longrightarrow N/(J\cap C_m)
 \quad\text{and}\quad N/(J\cap C_0)=N/J
\]
have finite $3$-group kernels. Pull back the first extension
along an already chosen section. Its quotient is the full
$G_{K_\infty}$; its finite kernel has a continuous action.
The preceding central-kernel induction applies verbatim:
finite-kernel extensions of profinite groups admit continuous
set-theoretic sections, and the elementary abelian step is
exactly continuous $H^2$. We may thus lift the chosen section
at each step. Their inverse limit is a continuous section of
$N\rtimes G_{K_\infty}$ and has the specified image in $N/J$.
This argument does not infer compatibility from unrelated
finite solutions.

The same absence of an inverse-limit defect can be seen
directly for coefficient lattices. The kernel of
$M/3^{m+1}\to M/3^m$ is finite $3$-primary. The dimension
bound makes its $H^2$ zero, so
\[
 H^1(G_{K_\infty},M/3^{m+1})
       \longrightarrow H^1(G_{K_\infty},M/3^m)
                     \quad\text{is surjective}.
\]
The first derived limit of these $H^1$ groups is therefore
zero, and all finite-coefficient $H^2$ groups are zero.
Jannsen's continuous-cohomology sequence proves the displayed
vanishing. No finiteness of these $H^1$ groups is required:
surjectivity itself gives the Mittag--Leffler property.
Thus the phantom obstructions below appear upon descent to
$\mathbb Q$, not already over the cyclotomic field.
\end{proof}

\begin{statement}{B.10}{Exact cyclotomic descent and its nonclosed image}{Proved}
Let $M$ be any finite free continuous $\mathbb Z_3[G_{\mathbb Q}]$-
module, let $H=G_{K_\infty}$, and put
\[
 \Delta=\{\pm1\}\subset\mathbb Z_3^\times,
 \quad\Gamma_1=1+3\mathbb Z_3,
 \quad\chi_3(\gamma)=4,
 \quad X=H^1_{\mathrm{cts}}(H,M)^\Delta.
\]
Give $X$ its $3$-adic inverse-limit topology. Then
\[
 H^2_{\mathrm{cts}}(G_{\mathbb Q},M)
       \simeq X/(\gamma-1)X,                                      \tag{B.10.1}
\]
where the denominator is the literal image, with no closure. The
kernel of reduction to all finite coefficient levels is
\[
 \frac{\overline{(\gamma-1)X}^{\,3\text{-adic}}}{(\gamma-1)X}
       \simeq\varprojlim{}^1_m H^1(G_{\mathbb Q},M/3^mM).           \tag{B.10.2}
\]
This includes the finite-free coefficient lattices of the class-two
and class-three obstruction sequences, for ramified inputs as well
as unramified ones.
\end{statement}
\begin{proof}
Write $M_m=M/3^mM$ and $X_m=H^1(H,M_m)^\Delta$. B.9 gives
$H^2(H,M_m)=0$. The exact sequences
$0\to M_1\to M_{m+1}\to M_m\to0$ show that the $H^1$ transition
maps are surjective. The $H^0(H,M_m)$ groups are finite, hence
Mittag--Leffler. The continuous-cohomology exact sequence of
\citet[Proposition 1.6 and Theorem 2.2]{BJannsen1988} consequently
gives $H^2_{\mathrm{cts}}(H,M)=0$ and
\[
 H^1_{\mathrm{cts}}(H,M)=\varprojlim_mH^1(H,M_m),\qquad
 X/3^mX=X_m.
\]
The second equality also follows from multiplication by $3^m$ and
the just-proved vanishing of continuous $H^2$. Taking $\Delta$
invariants is exact because two is a unit.

Use the \emph{inverse-system} Hochschild--Serre spectral sequence
\citep[Theorem 3.3]{BJannsen1988}. The $H^0$ row consists of finite
modules; its relevant higher continuous cohomology over $\Gamma_1$
vanishes. The $H^1$ row is the surjective system $X_m$; its inverse
limit is $X$. There is no $H^2$ row. This gives
$H^2_{\mathrm{cts}}(G_{\mathbb Q},M)=H^1_{\mathrm{cts}}(\Gamma_1,X)$.
It would not be valid to invoke an unrestricted continuous-cochain
Hochschild--Serre assertion without these inverse-limit checks.

Evaluation at $\gamma$ identifies continuous cocycles with $X$.
Indeed any $x\in X$ has an open stabilizer modulo $3^m$, and a
sufficiently long $3$-power norm of $x$ is zero modulo $3^m$.
Thus the formula on integral powers of $\gamma$ extends
continuously to $\mathbb Z_3$ for every $x$. Coboundaries evaluate
to $(\gamma-1)X$. This proves (B.10.1).

At finite level, ordinary Hochschild--Serre gives
\[
 H^2(G_{\mathbb Q},M_m)
   =X_m/(\gamma-1)X_m
   =X/\bigl((\gamma-1)X+3^mX\bigr).
\]
The kernel of the comparison map is therefore
$\bigcap_m((\gamma-1)X+3^mX)/(\gamma-1)X$, precisely the closure
quotient in (B.10.2). Jannsen's comparison exact sequence identifies
the same kernel with the displayed $\varprojlim^1$ term.

There is an explicit cochain form for every obstruction $o$, including
the cup obstruction $\eta\cup\beta$. Represent it by an extension
$E_o$ of $G_{\mathbb Q}$ by $M$, and split its restriction to $H$,
using $H^2_{\mathrm{cts}}(H,M)=0$. Denote the splitting by $r$.
Choose $\sigma$ above $\gamma$ and a lift $\widetilde\sigma$ in
$E_o$. The continuous $H$-cocycle
\[
 h\longmapsto
 \widetilde\sigma r(\sigma^{-1}h\sigma)
          \widetilde\sigma^{-1}r(h)^{-1}
\]
defines, after the $\Delta$ projection, an element of $X$.
Its class modulo $(\gamma-1)X$ is the image of $o$ under (B.10.1).
Changing the splitting over $H$ changes this element by
$(\gamma-1)y$. Thus every R10 or R12 obstruction has an exact
descent representative by this formula; vanishing at every finite
level means membership in the closure, whereas a continuous lift
requires membership in the image itself. Computing its numerical
coordinates still requires the actual action cocycle, and (B.10.1)
does not evaluate that missing cocycle.
\end{proof}

\begin{statement}{B.11}{The restricted-family calculation in R12}{Proved}
Let $\mathcal R=\{(z_q)_{q\ne3}\in\prod_q\mathbb Z_3:z_q\to0\}$,
and let $d$ act by $(dz)_q=(q-1)z_q$. Its coefficient comparison
sequence is
\[
 0\longrightarrow
 \frac{\{z\in\mathcal R:(q-1)\mid z_q\text{ for every }q\}}{d\mathcal R}
 \longrightarrow\mathcal R/d\mathcal R
 \longrightarrow\varprojlim_m\bigoplus_{q\ne3}
             \mathbb Z/(3^m,q-1)\longrightarrow0.                \tag{B.11.1}
\]
The left term can be nonzero. The tower
$\bigoplus_q(\mathbb Z/3^m)[q-1]$ is not Mittag--Leffler, and
supplies precisely the derived-limit defect in this presentation.
This verifies the algebra of R12; identifying a particular arithmetic
obstruction with a stated vector still requires its arithmetic residue
calculation.
\end{statement}
\begin{proof}
Put $A_m=\bigoplus_q\mathbb Z/3^m$,
$K_m=\ker(d:A_m\to A_m)$ and $I_m=dA_m$.
Both $A_m$ and $I_m$ have surjective transition maps. Moreover
$\varprojlim A_m=\mathcal R$ and
\[
 \varprojlim I_m=
 \{z\in\mathcal R:z_q\in(q-1)\mathbb Z_3\text{ for every }q\}.
\]
The derived-limit sequence of $0\to K_m\to A_m\to I_m\to0$
identifies $\varprojlim^1K_m=(\varprojlim I_m)/d\mathcal R$.
Apply the inverse limit also to
$0\to I_m\to A_m\to A_m/I_m\to0$; since $I_m$ is
Mittag--Leffler it gives the asserted surjection and (B.11.1).

For a single coordinate with $t=v_3(q-1)$, the image from level
$m\geq n$ to level $n$ in $K_n$ is
$3^{\max(m-t,0)}(\mathbb Z/3^n)$, interpreted as zero if the
exponent is at least $n$. For fixed $n$ these images would stabilize
uniformly if $t$ were bounded. They do not: primes with
$v_3(q-1)$ arbitrarily large exist by Dirichlet's theorem, and at
each later level a coordinate that survived an earlier level is
lost. Already at $n=1$ the descending supports are the primes with
$t\geq m$, so they never stabilize.

For an explicit nonzero kernel class choose distinct primes $q_j$
with $v_3(q_j-1)\to\infty$, and put $z_{q_j}=q_j-1$, all other
coordinates zero. Then $z\in\mathcal R$ and every coordinate is
divisible by $q-1$. Its only coordinatewise preimage under $d$ has
value one at every $q_j$, so does not belong to $\mathcal R$.
The class is nonzero globally and zero at every coefficient level
and every separate local coordinate. Under the arithmetic
identification $H^2_{\mathrm{cts}}=\mathcal R/d\mathcal R$, this is
exactly the form of the phantom term in R12 and of the nonclosed
image in B.10.
\end{proof}

\subsection{Resolved Selmer families and the remaining lifting questions}

\begin{statement}{B.12}{The exact Selmer classification of the R9 families}{Proved}
For every lift produced through the unramified arithmetic source
in LM6 E.16 and LM7 E.34, Selmerness is equivalent to membership
in one of the three global rational cusp packets. This holds for
all $n=\ell^r$, $\ell\in\{3,5\}$, for the full geometric
pro-$\ell$ group and every actual stable quotient retaining its
abelianization. The packets are the three copies of
$\mathbb Z_\ell\kappa_\ell(\ell)$ described in B.27.
Every other member fails the local point-or-cusp condition
already at $\ell$. In particular the explicitly noncuspidal
families of LM6 E.16, LM7 E.34 and LM7\_PUSH B.2 are non-Selmer.

In the full profinite metabelian families of LM6 E.31,
E.33 and E.34, exactly the zero parameter is Selmer; it is the
specified cusp section, by B.6. The least positive derived
length requested for the full profinite geometric quotients is
$k=1$ for both $n=3$ and $n=5$, by B.5 and B.26.
\end{statement}
\begin{proof}
The unramified classification is proved in B.27 for every
continuous cocycle with that hypothesis, independently of how
the lift was constructed. Its proof checks directly that the
named constructions use this arithmetic source. B.6 handles
the separate full profinite metabelian families. B.5 and B.26
prove the two least-positive-length assertions; neither is a
monotonicity assertion in derived length. These three results
give exactly the stated classification. No classification of
arbitrary sections ramified away from the coefficient prime
is inferred from B.27.
\end{proof}
The maximal-prosolvable quintic question is also settled by
A.50--A.52 and B.1, through the genus-two target with its
precise published Hilbert-modular input. This uses a different
target from the mod-five Legendre neighbourhood: its
$\mathrm{SL}_2(\mathbb F_5)$ quotient does not factor through
the maximal prosolvable geometric group. At seven, B.4 uses
$X_0(49)$ instead.

\begin{statement}{B.13}{The remaining arbitrary-input class-three and intrinsic quotient maps}{Open}
The full class-three image requested in prompt B.5 and the
intrinsic metabelian and marked images requested in B.6 are
not determined by this work. Every nonempty class-three
fibre is evaluated in B.28. The entire zero-abelian slice
has the necessary and sufficient local-order and decay law
proved in B.31. Its first local order is $m_2=3$ by B.36;
B.40 evaluates the complete local action class at two.
The remaining local orders are specified arithmetic data,
with $1\le m_5\le5$ the smallest unevaluated numerical
instance. The class-two law in prompt B.7 is proved
separately in B.23--B.25.

Prompt B.4 and Lead 3 are resolved by B.19--B.22 and
B.37--B.40, with a correction to the proposed method.
The rational coefficient is $-1$ relative to the specified
nonzero Ross cup edge. The symbol's real-regulator coordinate
is $-3/2$ in the distinct normalization of B.19.
The three-dimensional kernel in LM6 E.39 is an abelian-path
kernel, not the ambient action torsion.
The actual geometric constraints reduce the homogeneous
action ambiguity from $(\mathbb F_3)^4$ to $\mathbb Z/3$,
and B.40 selects its unique actual phase. In the intrinsic
finite test of B.39 the exact value is $19$ modulo $27$.
The prescribed inherited order-$3^{30}$ quotient kills the
entire residual line and cannot make this separation.
\end{statement}
The integral refinement is determined as a class over the
explicit rational Ross cup edge, by the actual path identities,
the integral commutator module, and the geometric extension
computed in B.40. This is not a numerical evaluation of an
infinite three-adic regulator. No such extra numerical datum
is required to identify the integral refinement requested
in B.4. The unresolved arbitrary-input maps in B.5--B.6
involve additional lifting data; they are not declared
settled merely because the determinant-extension class is
now identified. B.7 supplies a liftable infinite family,
and B.10 supplies an exact obstruction space, with these
scopes retained.

\begin{statement}{B.14}{The remaining metabelian and cross-prime colouring tasks}{Open}
The full profinite metabelian section sets in B.9, equality
of cusp colours at distinct ramified primes in B.10(a), and
the full pro-$3$ infinite-ramification equivalence in B.11
remain unproved here. B.29--B.30 refute the different assertion
that ramification at one prime forces a global rational cusp,
including its intrinsic metabelian and marked versions.
The explicit 44-coordinate phantom and the fibre reconciliation
are supplied in B.15--B.17 below.
\end{statement}
The first already requires the cubic at coefficient prime two,
where averaging by $1/2$ and the class-two BCH coordinates used
above are unavailable integrally. A proof for a single prime
does not give the full profinite metabelian group: mixed-prime
finite metabelian quotients must also be controlled. For cross-prime cusp
colourings the local descriptions identify $i(q)$ one
place at a time. The counterexample B.30 has one ramified
rational prime away from three and so does not decide equality
of labels at distinct primes. The smallest unresolved global
test is a section ramified at two good primes with different
local cusp labels. B.7 refutes the class-three half of B.11 but does not
refute its full pro-$3$ half: its central ramification with zero
abelian inertia cannot be the image of a ramified local cusp
in the full pro-$3$ group. The descent identities of B.10 cover
the obstruction groups in B.13 of the prompt; a numerical
rederivation of all of LM7 E.44 still requires the uncomputed
class-three arithmetic action.

\subsection{An explicit phantom in the forty-four coordinates}

\begin{statement}{B.15}{One nondecaying coordinate suffices}{Proved}
Use exactly the cusp order and quotient-wedge basis of LM7 E.37:
\[
 (f_1,f_2,e_{A_0},e_{A_1},e_{A_2},e_{B_0},e_{B_1},e_{B_2},
 e_{C_1},e_{C_2}),
\]
with basic wedges in lexicographic order, omitting
$e_{A_0}\wedge e_{A_2}$ in $M/L$. There is a continuous abelian
class whose unique rational corrections have precisely one
nonzero coordinate, the nineteenth, and every one of those
nonzero values is a $3$-adic unit. Its obstruction to class two
is nonzero continuously, zero locally everywhere, and zero at
every finite coefficient level.
\end{statement}
\begin{proof}
Put $P=A_0$ and $Q=B_0$, and let $\alpha_Q$ be the abelianized
class of the actual unit tangent at $Q$. It is unramified away
from $3$. For each $j\geq1$, let $p_j$ be the least unused prime
satisfying
\[
 p_j\equiv1+3^j\pmod{3^{j+1}},\qquad
                     3^{(p_j-1)/3}\not\equiv1\pmod{p_j}.          \tag{B.15.1}
\]
These conditions are soluble for every $j$. Indeed
$\mathbb Q(\zeta_{3^{j+1}})$ and
$\mathbb Q(\zeta_3,\sqrt[3]{3})$ intersect in
$\mathbb Q(\zeta_3)$. Otherwise the nonabelian $S_3$ splitting
field of $X^3-3$ would be contained in an abelian cyclotomic
extension. In the compositum prescribe cyclotomic Frobenius
$1+3^j$ and a nonidentity cubic Kummer Frobenius. They agree on
the intersection, and Chebotarev supplies infinitely many primes
in their conjugacy class. Excluding the finitely many previously
chosen primes proves existence in (B.15.1).

Set
\[
 \beta=\sum_{j\geq1}3^j\kappa_3(p_j),\qquad
 \alpha=\alpha_Q+\kappa_3(3)e_P+\beta e_Q
                   \in H^1(\mathbb Q,V).                         \tag{B.15.2}
\]
This is an actual continuous cocycle class: each cusp line is a
$G_{\mathbb Q}$-equivariant Tate line in $V$, and the indicated
sum converges uniformly. At a prime $\mathfrak p$ over $p_j$ of
$K=\mathbb Q(\zeta_3)$, put $q=p_j$ and choose the representative
of the framed cusp Kummer class in E37.10 to be
$k_{\mathfrak p,Q}=\alpha_Q(F_{\mathfrak p})$. The inertia is
$v_{\mathfrak p}=3^je_Q$. Substitution in that formula cancels
the cusp-path term literally and gives
\[
 D_{\mathfrak p}=3^j\kappa_3(3)(F_{\mathfrak p})
                       e_P\wedge e_Q.                            \tag{B.15.3}
\]
The $\beta(F_{\mathfrak p})e_Q$ term wedges to zero. On a pure
toric wedge the map $d_{\mathfrak p}=q^{-1}\Lambda^2A-1$ is
multiplication by $q-1$. Consequently the correction is exactly
\[
 c_{\mathfrak p}=
 -\frac{3^j\kappa_3(3)(F_{\mathfrak p})}{p_j-1}
                          [e_{A_0}\wedge e_{B_0}].                \tag{B.15.4}
\]
All other $43$ coordinates are zero. This wedge is retained in
the specified quotient basis, and is its nineteenth vector.
At every unlisted prime the inertia and the correction are zero.

Here are fully finite choices for the scalar in (B.15.4). Write
$u_j=(p_j-1)/3^j$, choose the least positive primitive root $g_j$
modulo $p_j$, and choose the local root-of-unity coordinate in
which $\zeta_{3^j}$ reduces to $g_j^{u_j}$. Let
\[
 0\leq b_j<3^j,\qquad
        (g_j^{u_j})^{b_j}\equiv3^{u_j}\pmod{p_j}.                 \tag{B.15.5}
\]
The unramified Kummer class of $3$ has Frobenius value $b_j$
modulo $(p_j-1)\mathbb Z_3$. This finite logarithm by itself
does not compute the full value of a fixed global Kummer
cocycle. We obtain an exact value by specifying the local
prolongations as follows. The restriction
$\kappa_3(3):\ker\chi_3\to\mathbb Z_3$ is surjective. Indeed
it is a continuous homomorphism with compact image; the image
is nonzero modulo three because $\sqrt[3]{3}$ is not contained
in the abelian cyclotomic tower. A closed subgroup of
$\mathbb Z_3$ containing a unit is the whole group.
For an initial Frobenius $F$, choose $h\in\ker\chi_3$ with
\[
 \kappa_3(3)(h)=
       \frac{\kappa_3(3)(F)-b_j}{p_j-1}\in\mathbb Z_3
\]
and replace the prolongation by its $h$-conjugate. The cocycle
law gives
$\kappa_3(3)(hFh^{-1})=b_j$ exactly. Its matrix on $V$ does
not change, since the image of $G_K$ on $V$ is abelian, and
$h$ fixes the root-of-unity coordinates. With these actual
prolongations, the nineteenth correction is the exact rational
$3$-adic number $-b_j/u_j$.
For arbitrary fixed prolongations it differs by $-3^jt_j$,
with $t_j\in\mathbb Z_3$, a family tending to zero. Thus the
nondecay conclusion is independent of this normalization.
The conjugate prime is obtained by complex conjugation, giving
the conjugate Tate coordinate; its integrality and unit valuation
are unchanged. Formula (B.15.1) says precisely that $3\nmid b_j$;
also $3\nmid u_j$. The correction is integral at every prime and
does not tend to zero.

For reference the executed first values are
\[
\begin{array}{c|r|r|r|r}
 j&p_j&u_j&b_j&c_{\mathfrak p,19}\\\hline
 1&13&4&1&-1/4\\
 2&37&4&8&-2\\
 3&109&4&25&-25/4\\
 4&811&10&1&-1/10\\
 5&5347&22&1&-1/22\\
 6&2917&4&697&-697/4\\
 7&87481&40&52&-13/10\\
 8&183709&28&6488&-1622/7.
\end{array}
\]
The ordering is by the exponent $j$, not by the magnitude of
$p_j$. The entire forty-four-coordinate vector is therefore
$(0^{18},-b_j/u_j,0^{25})$ in these local normalizations.

For completeness, the global obstruction statement follows
without assuming the full inherited E.37 calculation. Let
$M=\Lambda^2V$, let $L=\mathbb Z_3w$ be its primitive compact
determinant line, and let
$M_0=\ker(M\to\Lambda^2T_3E)$. The same finite-support
cohomology argument as in B.7 gives
$H^2_{\mathrm{cts}}(G_{K,S},M)=0$. The continuous localization
complex has restricted-family terms. At every good prime,
the kernel of $q^{-1}\Phi_M-1$ is $L$ and the image is in
$M_0$; on $M/L\otimes\mathbb Q_3$ this operator is invertible.
Expanding the class-two tame relation with the actual cusp
path at $Q$ gives exactly (B.15.3). A global continuous lift
would give restricted integral preimages, hence restricted
corrections in $M/L$. Formula (B.15.4) disproves this.
At each finite coefficient level only finitely many residues
survive, and each has an integral preimage, so that obstruction
vanishes. Local finite $H^1$ groups are Mittag--Leffler; thus
continuous local $H^2$ injects into the limit of its finite
coefficient groups, proving local vanishing also at $3$.
Restriction--corestriction transfers the claims from $K$ to
$\mathbb Q$, since its degree is two.

The file \texttt{b\_phantom44.py} executes the prime search,
the finite-field logarithms, the complete coordinate ordering,
and the exact residue equations. In particular the finite
precision does not stand in for a proof of rational recovery:
$-b_j/u_j$ is the exact correction in the prolongations just
specified, rather than a reconstruction of an unspecified
$3$-adic Frobenius value from a finite logarithm.
\end{proof}

\begin{statement}{B.16}{The phantom comparison and the missing free coordinate}{Proved}
The class of LM7\_PUSH A.4 maps to a phantom of the type in the
LM7 E.37 criterion. The two fibre ranks are compatible:
\[
\begin{aligned}
 H^1(G_{\mathbb Q,S},M_0)
    &\simeq\mathbb Z_3^{23}\oplus(\mathbb Z/3)^9
                            \oplus(\mathbb Z/9)^5,\\
 H^1(G_{\mathbb Q,S},M)
    &\simeq H^1(G_{\mathbb Q,S},M_0)\oplus\mathbb Z_3.
\end{aligned}                                                     \tag{B.16.1}
\]
The first fibre is over a specified cuspidally central section;
the second is over its abelianization. The extra coordinate is
the unramified compact-determinant Kummer coordinate, not an
extra phantom or a discrepancy in the finite torsion.
\end{statement}
\begin{proof}
The primitive $w$ has compact determinant $9$, so the
natural injection
\[
       \theta:M_0\lhook\joinrel\longrightarrow M/L
       \quad\hbox{has cokernel }\mathbb Z/9.                     \tag{B.16.2}
\]
Indeed the cokernel is
$M/(M_0+L)\simeq\mathbb Z_3/9\mathbb Z_3$, and the
intersection $M_0\cap L$ is zero. This integral finite-index
comparison is essential; rational identification alone would
not compare convergence.

Here is a direct check of the arithmetic construction used in
LM7\_PUSH A.4. Its pullback lattice is
$0\to M_0\to E_w\to\mathbb Z_3(1)\to0$ for the central
inertia line $c_O+c_P$. For the framed point $z_P$, tangent-line
evaluation gives the torus coordinates
\[
 3z_P=(3^{-2},-3^{-1},\lambda^{-1},\lambda^{-2},
                                  \zeta^2\lambda^{-3}),
       \qquad \lambda=1-\zeta.
\]
Take $\tau\in\ker\chi_3$ with $\kappa_3(3)(\tau)$ a unit
and $\sigma\in G_K$ with $\chi_3(\sigma)=4$. Such a $\tau$
exists because the nonnormal pure cubic field is not contained
in the abelian cyclotomic tower. For $g=[\sigma,\tau]$,
$\kappa_3(3)(g)=3\kappa_3(3)(\tau)$, while $g$ acts trivially
on $V$: its rational compact CM representation and its toric
representation are both abelian over $K$. In the non-$O$ cusp
basis used in that construction, the framed-point equality
therefore gives
\[
 3b_P(g)=\kappa_3(3)(g)v_0,\quad
 v_0=-2e_P-e_{Q_c}-\tfrac12(e_A+e_{\bar A})
                -(e_B+e_{\bar B})-\tfrac32(e_C+e_{\bar C}).
\]
The off-diagonal extension vector is
$u_g=-\kappa_3(3)(\tau)v_0\wedge e_P$, primitive because its
$e_P\wedge e_{Q_c}$ coefficient is a unit. Its action on
$M_0(-1)$ is the identity. Chebotarev applied modulo $3^j$
thus gives primes with $D_q\equiv0\pmod{3^j}$ and $u_q$
primitive. Choosing the least exponent $e_j$ for which
$3^{e_j}D_q^{-1}u_q$ is integral makes these corrections
primitive, while $e_j\geq j$ makes the input Kummer family
restricted. This proves directly the asserted nonzero phantom
in $H^2(\mathbb Q,M_0)$ by the continuous residue complex.

Under (B.16.2) its corrections are exactly the E.37 corrections,
up to a restricted family from the differing choices of cusp
inertia lift. They remain integral. They cannot tend to zero:
if $\theta(x)\in3^m(M/L)$, the cokernel killed by nine implies
$9x\in3^mM_0$. A primitive $x$ therefore cannot have
$\theta(x)\in27(M/L)$. Equivalently the kernel of
$H^2(\mathbb Q,M_0)\to H^2(\mathbb Q,M)$ is killed by nine,
whereas the nonclosed-image phantom subgroup is torsion-free.
Thus the obstruction does not disappear when one forgets the
specified cuspidally central lift and retains only its
abelianization. This is the required comparison with LM7\_PUSH
A.4, rather than an inference solely from two nonzero local tests.

For the fibres, use the integral exact sequence
\[
 0\longrightarrow M_0\longrightarrow M\longrightarrow
                       \mathbb Z_3(1)\longrightarrow0.
\]
There are no invariant vectors in the last module, and
$H^2(G_{\mathbb Q,S},M_0)=0$ by the Euler-characteristic and
devissage argument of B.7. Since
$H^1(G_{\mathbb Q,S},\mathbb Z_3(1))=
\mathbb Z_3\kappa_3(3)$, we obtain a short exact sequence
\[
 0\to H^1(G_{\mathbb Q,S},M_0)\to H^1(G_{\mathbb Q,S},M)
                         \to\mathbb Z_3\to0.
\]
It splits as a sequence of $\mathbb Z_3$-modules and preserves
all finite torsion. The explicit arithmetic matrices of E.37
give the presentation $(B_\tau-1)(1+B_c)$ on $M$ and its
restriction on $M_0$. Independent exact integer Smith
reduction in \texttt{b\_phantom44.py} finds in both cases seven
unit factors, nine factors of valuation one, and five of
valuation two; the free ranks are respectively $24$ and $23$.
This verifies (B.16.1) in the same integral lattice.
The arithmetic generators used in this presentation can also
be checked without a table. At the two primes over seven, take
the reductions $\zeta=2$ and $\zeta=4$. The Kummer coordinates
of $(\zeta,1-\zeta)$ on their Frobenii are respectively
$(2,0)$ and $(2,2)$ modulo three. They are linearly independent.
Thus these Frobenii generate the rank-two free pro-$3$
arithmetic source by the Frattini criterion; complex conjugation
interchanges them. Normalizing the cocycle at conjugation gives
precisely the quotient by $(B_\tau-1)M^c$ used in the Smith
calculation, and similarly for $M_0$.

Finally $M_0$ has no nonzero good-prime inertia eigenvectors,
so its global $H^1$ is already unramified away from $3$.
For $M$, the good-prime inertia line is precisely $L$, and the
cocycles $\sum_{p\ne3}a_p\kappa_3(p)w$ split its restricted
residue coordinates. Consequently the full abelian fibre is
$\mathcal R(\mathbb Z_3)\oplus H^1(G_{\mathbb Q,S},M)$,
as in LM7 E.37, whereas the fibre over a fixed cuspidally
central section is $H^1(G_{\mathbb Q,S},M_0)$, as in
LM7\_PUSH A.7--A.8. The map on the ramified determinant
coordinate is multiplication by nine, because the compact
projection of $w$ is nine. Thus the two descriptions are
compared through their actual coefficient maps.
\end{proof}

\begin{statement}{B.17}{The same failure of continuity in degree three}{Proved}
The central class-two family in LM6 E.38 is a degree-three
instance of the nonclosed-image phenomenon in B.11 and B.15.
In its actual integral lattice, its connecting class is nonzero
in continuous cohomology, zero at every finite coefficient
level, and zero over every completion. This assertion does not
require a value for the elliptic Ross component.
\end{statement}
\begin{proof}
We supply the two arithmetic calculations needed to justify
the named earlier example. Set $A=\operatorname{Lie}_3(V)$ and
choose the integral lift $\omega$ of $w$ with
$c\omega=-\omega$. Such a lift exists by averaging with $c$,
since two is a unit. Define the normalized genuine cocycle
\[
 \eta(g)=\chi_3(g)^{-1}g\omega-\omega
                              \in A(-1).
                                                               \tag{B.17.1}
\]
The class-two section $\beta w$ has connecting class
$[\eta\smile\beta]\in H^2_{\mathrm{cts}}(\mathbb Q,A)$:
this follows directly by lifting it to $\beta(g)\omega$ in the
abelian extension $\gamma_2N/\gamma_4N$ and taking its
continuous coboundary. This calculation uses the actual
extension; it is not a replacement by its associated graded.

First we compute an integral scalar detector. Put
$f=(X+Z)/(X+Y)$ and let $\ell_f:V\to\mathbb Z_3(1)$ be its
geometric Kummer map. Its boundary values are
$3(\delta_{B_0}-\delta_O)$ and its compact values are fixed
by vanishing on the transferred compact summand. Hence in the
basis of B.15,
\[
              \ell_f=(0,2,0,0,0,3,0,0,0,0).
\]
Compose the inclusion of the free degree-three Lie lattice in
$V^{\otimes3}$ with $\ell_f\otimes\det(p,p)$, where
$p:V\to T_3E$ is compactification, and denote the result by
$\Phi:A\to\mathbb Z_3(2)$. On the Hall basis
$[[e_i,e_j],e_k]$, $i>j$, $k\geq j$, its only nonzero entries
are $6,6,3$ at $(1,0,1),(1,0,5),(5,0,1)$. In particular
$\Psi=\Phi/3$ is an integral equivariant surjection, and after
twisting is a map $A(-1)\to\mathbb Z_3(1)$.

We prove its value on (B.17.1) geometrically. The $54$
automorphisms $\mathcal A=E[3]\rtimes\mu_6$ preserve $w$.
For $a\in\mathcal A$, choose a base-point path and let
$\alpha_a$ be the abelian path class of the image tangent
$a(s_O)$. Applying the identity
$\sigma_a\rho_g\sigma_a^{-1}=\operatorname{Int}(k_a(g))\rho_g$
to $\omega$, and taking degree three, gives
\[
       (a-1)[\eta]=[\alpha_a,w].
\]
The ambiguity from $\sigma_a\omega-\omega$ is exactly a
coboundary. Thus, with the rational averaging idempotent
$e_{\mathcal A}$ and
$\overline\alpha=54^{-1}\sum_a\alpha_a$,
\[
          (1-e_{\mathcal A})[\eta]=-[\overline\alpha,w].
                                                               \tag{B.17.2}
\]
For each nonorigin cusp $d$, the tangent-line unit
$f_d=L_d/(X+Y)$ has divisor $3(d-O)$. Every line occurs six
times among the monomial automorphisms, so
$\prod_a a^*f_d$ is a root of unity. Pairing the mean path
class with $f_d$ therefore gives
$-\kappa_3(\operatorname{lt}_O f_d)$ in rational cohomology.
At the specified unit tangent $q=Z/X$,
\[
 Y/X=-(1+q^3)^{1/3},\qquad
          \operatorname{lt}_O(f_d)=-3L_d(1,-1,0).
\]
For the six $A_i,B_i$ cusps its Kummer value is
$\kappa_3(3)$; for $C_1,C_2$ it is
$\tfrac32\kappa_3(3)$, using
$(1-\zeta)^2=-3\zeta$. Since the eight boundary forms are
three times the individual nonorigin coordinates, this proves
\[
 \overline\alpha=-\left(\tfrac13v_6+\tfrac12v_2\right)
                       \kappa_3(3),\quad
 v_6=\sum_i(e_{A_i}+e_{B_i}),\quad v_2=e_{C_1}+e_{C_2}.
                                                               \tag{B.17.3}
\]
No regulator theorem enters this calculation.

The map $\Phi$ is supported on the multidegree $H^2W$ for
the rational compact and boundary decomposition $V=H\oplus W$.
Translations by $E[3]$ act trivially on $H$ and have no
invariants on the augmentation representation $W$ of the nine
cusps. Consequently $\Phi e_{\mathcal A}=0$. Tensor expansion
also gives $\Phi([z,w])=18\ell_f(z)$ because $p(w)=9$.
As $\ell_f(v_6/3+v_2/2)=1$, (B.17.2)--(B.17.3) imply
$\Psi_*[\eta]=6\kappa_3(3)$ rationally. The integral group
$H^1(G_{\mathbb Q,\{3,\infty\}},\mathbb Z_3(1))$
is torsion-free and equals $\mathbb Z_3\kappa_3(3)$, so the
equality holds integrally. Both cocycles vanish at $c$ when
the Kummer roots of three are positive real roots. A possible
Tate coboundary has value $-2b$ at $c$, forcing $b=0$. We have
therefore proved the literal cochain identity
\[
                         \Psi\eta=6\kappa_3(3).                 \tag{B.17.4}
\]

Second we verify the uniform inverse bound. The actual
Frobenius matrix at seven on $V$ is the matrix $T$ used in
\texttt{b\_phantom44.py}; its compact block is
$I+3\left(\begin{smallmatrix}0&-1\\1&-1\end{smallmatrix}\right)$
and its toric block is $7I_8$. Put
\[
       R=7^{-1}\operatorname{Lie}_3(T),\qquad D=R^3-I.
\]
The independent program \texttt{b\_class3\_phantom.py}
constructs the complete $330$-element Hall basis and every
entry of $D$; it uses only the integral $10$-dimensional
matrix, not a matrix file from the previous return. Unimodular
reduction modulo $3^{12}$ gives the elementary-divisor
valuations
\[
                   1^{132},\ 2^{82},\ 3^{52},\
                   4^{49},\ 5^{14},\ 6^1.                     \tag{B.17.5}
\]
Every pivot is below the precision. To check the determinant
independently, write $\pi=1+3\zeta$, so that
$\pi+\bar\pi=-1$ and $\pi\bar\pi=7$. The eigenvalue blocks
give the five contributions
\[
\begin{gathered}
 v_3N(\pi^3-1),\quad 8v_3N(\pi^6-1),\quad16v_3(7^3-1),\\
 64v_3N((7\pi)^3-1),\quad168v_3(7^6-1),
\end{gathered}
\]
namely $4,32,32,320,336$, whose sum is $724$ and agrees
with (B.17.5). The row operations and this determinant check
certify the integral bound $3^6D^{-1}$. Moreover
$D\in3M_{330}(\mathbb Z_3)$, and
\[
 (I+B)^3-I=3B(I+B+B^2/3),\qquad
 I+B+B^2/3\in I+3M_{330}(\mathbb Z_3)
\]
for every $B\in3M_{330}(\mathbb Z_3)$. Therefore
\[
\begin{gathered}
 D_n=R^{3^{n+1}}-I=3^nDU_n,\qquad
 U_n\in\operatorname{GL}_{330}(\mathbb Z_3),\\
                 3^{n+6}D_n^{-1}\text{ is integral}.            \tag{B.17.6}
\end{gathered}
\]

We now specify the primes, retaining the actual normalization
in LM6 E.38. Start with a Frobenius $\tau_0$ at
$(7,\zeta-2)$ and put $\tau'_0=c\tau_0c$. Its positive-root
Kummer value $b=\kappa_3(3)(\tau_0)$ satisfies $b\equiv1\pmod3$
because $3^2\equiv2\pmod7$. For
$r_0=\tau'_0\tau_0^{-1}$ and
$h=r_0^{(b^{-1}-1)/12}$, the exponent is integral and the
cocycle law gives
$\kappa_3(3)(h\tau_0h^{-1})=1$. Take
$\tau=h\tau_0h^{-1}$ and $\tau'=c\tau c$. Their matrices
on $V$ remain $T$ and its conjugate, since the representation
over $K$ is abelian. Put $r=[\tau,\tau']$ and
$g_n=r\tau^{3^{n+1}}$ for $n\geq1$. Directly,
\[
 \rho_V(r)=1,\quad\chi_3(r)=1,\quad\kappa_3(3)(r)=-12,
\]
and hence
\[
\begin{gathered}
 \rho_V(g_n)=T^{3^{n+1}},\quad
 v_3(\chi_3(g_n)-1)=n+2,\\
                 v_3(\kappa_3(3)(g_n))=1.                      \tag{B.17.7}
\end{gathered}
\]
For the last equality use
$\kappa_3(3)(g_n)=-12+(7^{3^{n+1}}-1)/6$.

In the finite image of the actual representations
$\rho_V$, $\chi_3$, and $(\kappa_3(3),\chi_3)$ modulo
$3^{n+8}$, let $p_n$ be the least prime larger than $p_{n-1}$,
with $p_0=3$, whose Frobenius class equals the class of $g_n$.
Chebotarev proves existence. These primes split in $K$.
For suitable prolongations,
\[
\begin{gathered}
 D_{p_n}-D_n\in3^{n+8}M_{330}(\mathbb Z_3),\\
 v_3(p_n-1)=n+2,\qquad v_3\kappa_3(3)(F_{p_n})=1.
\end{gathered}
\]
The Neumann-series argument applied to
$D_n^{-1}(D_{p_n}-D_n)\in9M_{330}(\mathbb Z_3)$ preserves
the inverse bound in (B.17.6).

Thus the continuous cocycle
$\beta=\sum_{n\geq1}3^{n+6}\kappa_3(p_n)$ has integral
degree-three corrections
\[
             x_n=-3^{n+6}D_{p_n}^{-1}\eta(F_{p_n}).              \tag{B.17.8}
\]
Their images under the integral detector satisfy
\[
 \Psi x_n=-\frac{3^{n+6}\,6\kappa_3(3)(F_{p_n})}{p_n-1},
                         \qquad v_3(\Psi x_n)=6.               \tag{B.17.9}
\]
Consequently they do not tend to zero. The degree-three
continuous localization complex of B.7 has precisely this
restricted-family condition: (B.17.8) are the unique rational
preimages, so their nondecay proves a nonzero continuous
obstruction. For every finite partial sum they are integral
and finitely supported, and the obstruction vanishes. Modulo
$3^m$ the full obstruction equals that of a finite partial
sum, proving zero at every coefficient level. Finite local
$H^1$ groups are Mittag--Leffler, so continuous local $H^2$
injects into their coefficient limit; this proves local
vanishing, including at three. Averaging handles the real
place.

This is exactly the difference between membership in the
closure of a residue image and membership in that image.
It is the same cohomological mechanism as B.15, with
$\Lambda^2V$ replaced by $\operatorname{Lie}_3(V)$ and the
quadratic tame equation replaced by the connecting class
$\eta\smile\beta$. It does not identify the two classes in
different coefficient groups. The stated arithmetic matrices,
the scalar detector, and the determinant certificate have
all been checked independently; no elliptic regulator or
uncomputed integral Ross coordinate was used.
\end{proof}

\begin{statement}{B.18}{The infinite compact cubic families fail at three}{Proved}
Let $U_3$ be the plus Fermat cubic with its nine cusps removed,
and let $s$ be a section in the LM6 E.16 or LM7 E.34
construction whose compact abelian class
$a\in H^1(G_{\mathbb Q,\{3,\infty\}},T_3E)$ has infinite
additive order. Every full pro-$3$ lift of that class fails the
Selmer condition at three. The same is true in each quotient
which retains its compact Tate-module image.
This argument alone treats the infinite-compact-order
subfamily. The complete classification of all unramified
lifts, including torsion compact classes and the higher
$3$-power and $5$-power levels, is supplied by B.27 below.
\end{statement}
\begin{proof}
At every prime $q\ne3$, the elliptic curve $E=X_3$ has good
reduction and the class $a$ is unramified. For good reduction
and coefficient prime different from the residue prime,
the unramified subgroup of $H^1(\mathbb Q_q,T_3E)$ is exactly
the image of $E(\mathbb Q_q)^{\wedge_3}$ under Kummer.
Indeed reduction identifies the latter completion with
$E(\mathbb F_q)^{\wedge_3}$, since the formal group is pro-$q$.
The Kummer sequence over the finite residue field, together
with surjectivity of Frobenius minus one on
$E(\overline{\mathbb F}_q)$, identifies this group with
$T_3E/(F_q-1)T_3E$, the unramified cohomology group.
The real place has zero positive-degree $3$-primary
cohomology.

If $s$ were locally a point or cusp at three, functoriality
under compactification and the Abel--Jacobi map would put
$a|_{G_{\mathbb Q_3}}$ in the Kummer image of
$E(\mathbb Q_3)^{\wedge_3}$. Its reductions modulo $3^m$
would therefore lie in all the global $3^m$-Selmer groups.
The finite Kummer exact sequences
\[
 0\longrightarrow E(\mathbb Q)/3^mE(\mathbb Q)
 \longrightarrow\operatorname{Sel}_{3^m}(E/\mathbb Q)
 \longrightarrow\Sha(E/\mathbb Q)[3^m]\longrightarrow0
\]
are \citet[I, Proposition 6.4]{BMilne2006}. Taking inverse
limits is exact on the left terms, which are finite and
Mittag--Leffler, and gives
\[
 0\longrightarrow E(\mathbb Q)^{\wedge_3}
 \longrightarrow\varprojlim_m\operatorname{Sel}_{3^m}(E/\mathbb Q)
 \longrightarrow T_3\Sha(E/\mathbb Q)\longrightarrow0.
\]
Part A proves that $E(\mathbb Q)$ and $\Sha(E/\mathbb Q)$
are finite. Hence the last term is zero and the middle
group is finite. The usual continuous-cohomology comparison
in degree one has no extra derived-limit term: the invariant
groups $H^0(\mathbb Q,E[3^m])$ are finite and
Mittag--Leffler. Thus $a$ itself would have finite additive
order, a contradiction.
The argument uses only its compact image. Every section
which maps to a local point or cusp in any of the stated
quotients has a local compact Kummer image, so the same
contradiction applies there.
\end{proof}

\subsection{What the Ross normalization does determine}

\begin{statement}{B.19}{The exact Hodge coordinate of the Ross symbol}{Known}
On $F_3:u^3+v^3=1$ put $e_R=\{1-u,1-v\}$ in rational
motivic cohomology and let
\[
 \varphi:F_3\xrightarrow{\sim}E_{27},\qquad
 (u,v)\longmapsto
 \left(\frac{3v}{1-u},\frac{9(1+u)}{2(1-u)}\right),\qquad
 E_{27}:y^2=x^3-\frac{27}{4}.
\]
Set $e_{27}=\varphi_*e_R$. Let $\omega$ be the real
holomorphic differential normalized by
\[
 \frac{1}{2\pi i}\int_{E_{27}(\mathbb C)}
              \omega\wedge\overline\omega=-1,
 \qquad \Omega_{\mathbb R}
          =\int_{E_{27}(\mathbb R)^0}\omega>0.
\]
Then the real regulator has the exact coordinate
\[
 r_{\mathcal D}(e_{27})
   =-\frac32\,L'(E_{27},0)\Omega_{\mathbb R}
                         (\omega-\overline\omega).              \tag{B.19.1}
\]
Thus the coordinate is $-3/2$ in this specified $L$-value
normalization; the integral tame multiple $6e_R$ has
coordinate $-9$. This is \citet[Theorem 3.3]{BIto2018}, with
the conventions of his Section 3.1. It is a value of the
Ross symbol, not an evaluation of the integral arithmetic
action at a Frobenius.
\end{statement}
\begin{proof}
The inverse of $\varphi$ on the displayed affine charts is
\[
          u=\frac{2y-9}{2y+9},\qquad
          v=\frac{6x}{2y+9}.
\]
Substitution checks both inverse compositions and the curve
equation. A birational map between these smooth proper
curves extends to an isomorphism, so no degree or isogeny
factor is omitted. These identities are executed in
\texttt{b\_ross\_hodge.py}.
The symbol and its pushforward are exactly the elements
defined in Ito's Section 3.1; the area and orientation
normalizations displayed above are his normalizations.
His Theorem 3.3 therefore gives (B.19.1) without a change
of differential or regulator convention. Multiplying by six
gives $-9$ by linearity. The reason for retaining this factor
is that \citet[Theorem 1]{BRoss1994} supplies the tame integral
multiple $2N e_R$, which is $6e_R$ when $N=3$.
\end{proof}

This calculation replaces the vague instruction to read a
rational number from a Hodge realization by one fully
normalized real-regulator number. It does not identify a
$\mathbb Z_3$-basis of the arithmetic cohomology lattice.
The distinction is visible already in the primary literature:
\citet[Section 1]{BNemoto2025} records regulator equalities
for conductor $27$, while proving motivic equalities for his
different conductor-$36$ and conductor-$64$ examples.
Neither a real-regulator equality nor tensoring with
$\mathbb Q_3$ retains integral torsion. B.21 distinguishes the
three-dimensional path ambiguity from the larger action-cohomology
torsion group.
B.37 proves arithmetic nonvanishing by an independent
integral local calculation; B.38 and B.40 then determine the
integral refinement by actual geometric data. These results
preserve the distinction between the real normalization and
the arithmetic class.

\begin{statement}{B.20}{The rational unit-cup comparison, with its sign}{Proved}
For the cubic unit tangent and lattice conventions of B.17,
write $x=-X/Z$, $y=-Y/Z$, and put
$f=1-x$, $h=y-1$. On the compact and boundary splitting
$V_{\mathbb Q_3}=H\oplus W$, let
\[
 b=\ell_f\wedge\ell_h,\qquad
 J:H\otimes\Lambda^2W\longrightarrow\operatorname{Lie}_3(V),
       \quad J(a\otimes(r\wedge s))=[a,[r,s]].
\]
The invariant subspace $D=(\operatorname{im}J)^{\mathcal A}$
has dimension two, and
$\Phi=(1_H\otimes b)J^{-1}:D\to H(2)$ is a
$G_{\mathbb Q}$-equivariant isomorphism. If
$\rho=\kappa(f)\smile\kappa(h)$ and the mixed edge convention
is the one below, then
\[
       \Phi_*[\eta]_{\mathrm{ell}}
       =-\iota_{w*}\operatorname{edge}_H(\rho)
                          \quad\text{in }H^1(\Gamma,H(1)),
 \qquad
       \iota_w(\lambda)=u\lambda(v)-v\lambda(u).                 \tag{B.20.1}
\]
Here $\Gamma=G_{\mathbb Q(\zeta_3),\{3\}}(3)\rtimes
\langle c\rangle$ is the actual arithmetic source, and
$w=u\wedge v$. The coefficient relative to the specified
Ross cup-edge class is exactly $-1$, or $-1/6$ relative to
the edge of its tame multiple $6e_R$. This identity alone
does not establish arithmetic nonvanishing; B.37 proves
that the edge is a nonzero generator of its one-dimensional
rational arithmetic space.
\end{statement}
\begin{proof}
The tangent-line divisors are
\[
 \operatorname{div}(f)=3B_0-C_0-C_1-C_2,\qquad
 \operatorname{div}(h)=3A_0-C_0-C_1-C_2.
\]
They define the two boundary characters literally. These
characters vanish on $H$, because they are equivariant and
the compact and Tate weights at a good Frobenius differ.
In the unit parameter $q=Z/X$ at $O=[1:-1:0]$,
\[
 f=q^{-1}(1+q),\qquad
 h=q^{-1}\bigl((1+q^3)^{1/3}-q\bigr);
\]
their leading coefficients are one. Thus their arithmetic
Kummer cocycles vanish on the chosen tangential section.

In rational class-two Malcev coordinates the group law and
the actual arithmetic action have the forms
\[
\begin{aligned}
 (v,z)(v',z')&=(v+v',z+z'+\tfrac12v\wedge v'),\\
 g(v,z)&=(A_gv,(\Lambda^2A_g)z+B_gv),\qquad
 C_g=B_gA_g^{-1}.
\end{aligned}
\]
Composition of actions proves $C_{gg'}=C_g+gC_{g'}$.
Expanding $g[u,v]-\chi_3(g)[u,v]$ and retaining one
compact and two boundary letters gives
\[
 \eta_{HW^2}(g)=
 J\bigl(u\otimes C_g^{WW}(v)-v\otimes C_g^{WW}(u)\bigr).
                                                               \tag{B.20.2}
\]
This uses the determinant of $A_g|H$, which is $\chi_3(g)$.
Coordinate changes add the corresponding coboundaries.
The map $J$ is injective: the terms whose first letter lies
in $H$ recover the tensor $a\otimes(r\wedge s)$.

For elements $e=(n,g)$ and $e'=(m,g')$ of
$N\rtimes\Gamma$, write their geometric coordinates as
$(v_n,z_n)$ and $(v_m,z_m)$. The antisymmetrized cup is
$\omega(e,e')=\tfrac12 b(v_n\wedge A_gv_m)$.
Its difference from the ordered cup is a coboundary.
The continuous cochain $a(e)=-b(z_n)$ satisfies
\[
 \delta a(e,e')=\omega(e,e')+bB_g(v_m),
 \qquad
 \omega-\delta a=-(bC_g)(A_gv_m).                               \tag{B.20.3}
\]
The mixed cocycle on the right defines the edge as
$[-bC_g]$; in particular this computation fixes its sign.
The free pro-$3$ groups $N$ and the rank-two arithmetic
source have zero cohomology above degree one for these
rational coefficients. This follows first at finite
coefficients, then from finiteness of their finite
cohomology and inverse limits, followed by rationalization.
Averaging by $c$ is exact. The only Hochschild--Serre term
of total degree two is therefore
$H^1(\Gamma,\operatorname{Hom}(V,\mathbb Q_3(2)))$.
Restriction to $H$ gives the asserted compact edge.
Equations (B.20.2)--(B.20.3) now prove the signed identity
before taking invariants.

We check the invariant normalization independently. In a
rational compact CM basis put
$R=\left(\begin{smallmatrix}0&-1\\1&-1\end{smallmatrix}\right)$.
The automorphism consisting of diagonal exponents $(a,b,0)$
and a coordinate permutation $\pi$ acts on $H$ by
$\operatorname{sgn}(\pi)R^{a+b}$, as follows by pulling
back the holomorphic differential. On $W$ it permutes the
nine cusp vectors modulo their sum.
The program \texttt{b\_ross\_hodge.py} constructs all these
permutations in exact projective root-of-unity coordinates,
then their $56$-dimensional tensor-wedge actions. If $P$
is their sum, it verifies $P^2=54P$ and
$\operatorname{tr}P=108$. Thus $P/54$ has rank two.
Its two averaged columns
\[
 d_0=e_{\mathcal A}(a_0\otimes A_0\wedge B_0),\qquad
 d_1=e_{\mathcal A}(a_0\otimes A_0\wedge B_1)
\]
have contraction matrix
\[
          (1_H\otimes b)(d_0,d_1)
          =\begin{pmatrix}-1/3&1/3\\0&1/3\end{pmatrix},
                  \qquad\det=-1/9.                            \tag{B.20.4}
\]
It follows that the stated contraction is an isomorphism.
All these maps commute with Galois, which permutes the
automorphisms in the average. Finally B.17.2 shows that the
failure of $[\eta]$ to be invariant is a bracket
$[\alpha_a,w]$, with at least two compact letters.
Its $HW^2$ projection is zero. The elliptic cohomology
class is therefore unchanged by the average, proving
(B.20.1) on the specified invariant component.
Since $-1$ has compatible rational odd-power roots,
$\kappa_3(-1)=0$; replacing $h$ by $1-y$ identifies the
ordered cup with the étale symbol attached to $e_R$.
\end{proof}

The identity (B.20.1) compares two explicitly constructed
classes and fixes their rational scalar. B.37 proves that
the specified edge is a nonzero arithmetic generator, and
derives the dimension-one assertion directly at three.
B.38 and B.40 determine the integral refinement, culminating
in the exact intrinsic finite value $19$ modulo $27$.
Equation (B.19.1) remains the separate real-regulator
normalization; it is not a numerical three-adic value.

The theorem of \citet[Theorem 1.1.5]{BKings2001}
explicitly assumes $p\ne2,3$ and good reduction at $p$.
The all-prime Iwasawa theorem of
\citet[Theorem 5.2]{BJohnsonLeungKings2011}
does not by itself identify this specified Ross symbol
with an integral Frobenius cocycle. Neither theorem is
used outside its hypotheses to obtain the refinement.
Instead B.21--B.22 separate the path kernel and the
inherited-test kernel, B.38 gives the exact residual
cyclic ambiguity, and B.40 computes the actual geometric
extension that selects its phase.

\begin{statement}{B.21}{The three path ambiguities and the larger action lattice}{Proved}
The group $(\mathbb F_3)^3$ in LM6 E.39 is precisely the
kernel of the comparison of \emph{abelian path classes}
\[
 H^1(G_{\mathbb Q},V)\longrightarrow
 H^1(G_{\mathbb Q},T_3E\oplus\mathbb Z_3(1)^8),
                                                               \tag{B.21.1}
\]
where Galois permutes the eight unit coordinates. It is
also the kernel for the unramified arithmetic source.
It is not the torsion subgroup of the coefficient group
containing $\eta$. In the latter group one has instead
\[
\begin{aligned}
 H^1(\Gamma,\operatorname{Lie}_3(V)(-1))
  \simeq{}&\mathbb Z_3^{164}\oplus(\mathbb Z/3)^{59}\\
           &{}\oplus(\mathbb Z/9)^{27}
                    \oplus(\mathbb Z/27)^7.
\end{aligned}
                                                               \tag{B.21.2}
\]
An identification of the integral refinement of $\eta$
with only the three classes in (B.21.1) therefore requires
an additional theorem about the actual geometric extension;
LM6 E.39 is not that theorem. B.38 and B.40 provide the
actual action-class analysis and integral selection, without
identifying this path kernel with the action torsion.
\end{statement}
\begin{proof}
Write $U$ for the $8$-by-$2$ gluing matrix in
\texttt{b\_phantom44.py}. The normalized tangent-line unit
$f_d=L_d/L_O$ has boundary character $3e_d^*$ and vanishes
on the compact columns $\binom{3I_2}{U}$. Thus its full
character is $(-U_d,3e_d^*)$. The actual equivariant
comparison and quotient are
\[
 j=\begin{pmatrix}I_2&0\\-U&3I_8\end{pmatrix},
 \qquad
 (h,b)\longmapsto b+Uh\pmod3.
\]
Integral row reduction gives cokernel $(\mathbb Z/3)^8$.
The quotient action is trivial over $K=\mathbb Q(\zeta_3)$
and is minus the permutation
$(A_1\,A_2)(B_1\,B_2)(C_1\,C_2)$ over $\mathbb Q$.
The two fixed permutation coordinates $A_0,B_0$ contribute
no fixed vectors because the action there is minus one.
The invariant subspace therefore has basis
\[
 e_{A_1}-e_{A_2},\qquad e_{B_1}-e_{B_2},\qquad
 e_{C_1}-e_{C_2}.
\]
There are no invariant vectors in $T_3E$ or the unit
Tate module, as is seen at a good Frobenius over $K$.
The exact continuous cohomology sequence consequently
identifies this invariant quotient with the kernel
in (B.21.1). The same argument over $K$ gives
$(\mathbb Z/3)^8$. If $t_d$ is a boundary vector, the
connecting cocycle is
$g(t_d/3)-t_d/3$; the three displayed differences have
value zero at $c$ and value $2t_d$ at the normalized
Frobenius $\tau$ above seven. The matrix identities, all
ten Smith factors, and the three invariant vectors are
independently checked in \texttt{b\_phantom44.py}.

For (B.21.2), put $T_A=\operatorname{Lie}_3(T)$ and
$C_A=\operatorname{Lie}_3(C)$. A cocycle on the rank-two
free pro-$3$ arithmetic source can be assigned arbitrarily
at $\tau,\tau'=c\tau c$. Averaging normalizes its value
at $c$ to zero and forces the second value to be the
$c$-translate of the first. On $A(-1)$ the matrices are
$T_A/7$ and $-C_A$. Normalized coboundaries therefore give
the exact presentation
\[
 H^1(\Gamma,A(-1))
       =A\big/\bigl((T_A-7I)(I-C_A)A\bigr).                      \tag{B.21.3}
\]
The factor two in the invariant projector and the factor
seven in the twisted action are units, so this expression
uses the integral image exactly.
The trace of $c$ on $V$ is $-2$, hence its trace on
the third Lie power is $((-2)^3-(-2))/3=-2$.
Thus $I-C_A$ has rank $166$. The other factor is
invertible rationally by the good-Frobenius weights.
The independent Hall-basis program
\texttt{b\_class3\_phantom.py} reduces the actual matrix
in (B.21.3) and obtains $166$ nonzero Smith factors,
with valuations
$0^{73},1^{59},2^{27},3^7$. All occur below the working
precision $3^{12}$, and the exact rank calculation
proves that none are omitted. This gives (B.21.2).
\end{proof}

\begin{statement}{B.22}{Completeness of the fixed finite Frobenius test}{Refuted}
The images at all Frobenius elements in the order-$3^{30}$
quotient of LM7 E.38 do not determine a general integral
class in $H^1(\Gamma,A(-1))$, even after its rational image
is fixed. The induced map on that coefficient group has
an infinite kernel, and its kernel on the torsion subgroup
has order at least $3^{128}$. In particular that finite
test is not a complete invariant of the arithmetic
extension datum containing $\eta$.
\end{statement}
\begin{proof}
Denote the finite quotient by $Q=H_6/\gamma_4H_6$ and put
$A_Q=\gamma_3Q$. Its central coefficient map is the
continuous equivariant quotient
$k:A=\gamma_3N/\gamma_4N\to A_Q$. The quotient is built
inside
\[
 R=\mathbb Z_3[[T-1]]\rtimes\mathbb Z_3,\qquad
 g(a,v(T))=(\chi_3(g)a,\chi_3(g)v(T^{\chi_3(g)})).
                                                               \tag{B.22.1}
\]
For completeness this action is exact: in the tripod
action, $x\mapsto x^\chi$ and
$y\mapsto f_g^{-1}y^\chi f_g$ with $f_g\in F'$.
The image of $f_g$ lies in the abelian normal module
$I_y/[I_y,I_y]$, so its conjugation of $y$ disappears.
The resulting generator formulas are (B.22.1).
Changing the arithmetic origin from the tangent at $P$
to that at $O$ changes the action by an inner automorphism,
which acts trivially on the central group $A_Q$.
Thus $G_{\mathbb Q}$ acts on $A_Q$ through $\chi_3$
in either convention.

We first give an explicit infinite-order kernel element.
On $M=H(1)$ in the integral basis $(u,v)$, the matrices
at $\tau$ and $c$ are
\[
 T_M=7\begin{pmatrix}1&-3\\3&-2\end{pmatrix},
 \qquad
 C_M=-\begin{pmatrix}1&-1\\0&-1\end{pmatrix}.
\]
The invariant lattice for $c$ is generated by $(1,2)$,
and
\[
 (T_M-I)(1,2)^{\mathsf t}=-9(4,1)^{\mathsf t}.
\]
Hence
$H^1(\Gamma,M)=\mathbb Z_3^2/\mathbb Z_3\,9(4,1)$.
The continuous normalized cocycle $\xi$ given by
\[
 \xi(c)=0,\qquad \xi(\tau)=(1,0),\qquad
 \xi(\tau')=C_M(1,0)=(-1,0)
\]
exists by the actual free arithmetic source. It has
infinite order: the integral row $(1,-4)$ annihilates
all coboundaries and takes value one on it.

B.20's exact projector supplies an integral equivariant
injection $i:M\hookrightarrow A(-1)$ with
$\Phi i=18I$. Explicitly, if $P$ is the integral sum
of the $54$ tensor-wedge actions and $s_0,s_1$ are the
two seed columns in (B.20.4), take
\[
 i(e_1)=-JP s_0,\qquad i(e_2)=JP(s_0+s_1).
                                                               \tag{B.22.2}
\]
These formulas are integral, and multiplication by the
contraction matrix verifies $\Phi i=18I$.
Let $t=[c,\tau]$. Its cyclotomic character is one, so
it acts trivially on $A_Q(-1)$. On $M$ one computes
\[
 \rho_M(t)-I=\frac37
                  \begin{pmatrix}-4&-1\\1&-5\end{pmatrix},
 \qquad
 9(\rho_M(t)-I)^{-1}
                  =\begin{pmatrix}-5&1\\-1&-4\end{pmatrix}.
                                                               \tag{B.22.3}
\]
In particular its Smith factors over $\mathbb Z_3$
are $3$ and $9$. Equivariance implies
$ki(\rho_M(t)-I)M=0$, hence $9ki(M)=0$.
The two explicit cocycles $0$ and $9i\xi$ therefore
have identical images in $A_Q(-1)$ at every Galois
element, while their cohomology classes are distinct:
$\Phi_*(9i_*[\xi])=162[\xi]\ne0$ rationally.
All matrices in (B.22.2)--(B.22.3) and this detector
are checked in \texttt{b\_ross\_hodge.py}.
These are continuous arithmetic module-extension
classes. Their role is to disprove injectivity of the
finite test; they are not asserted to be two geometric
Galois actions of the same curve.

There is also a large loss within any fixed rational
class. We compute the size of the finite target.
Write $U=T-1$, $\delta=T^3-1$ and
$J=(3,U)$ in $(\mathbb Z/729)[U]/(U^5)$.
Then $A_Q=\delta^2J/\delta^3J$. In coordinates
$(U^2,U^3,U^4)$, the two lattices modulo $729$ are
generated by the columns
\[
 G_3=\begin{pmatrix}27&0&0\\54&9&0\\45&18&9\end{pmatrix},
 \qquad
 G_4=\begin{pmatrix}0&0\\81&0\\243&27\end{pmatrix}.
\]
Adjoin $729I_3$ to each, take their integral Hermite
bases $L_3,L_4$, and reduce $L_3^{-1}L_4$ to Smith
form. The exact result is $(3,9,27)$, independently
executed in \texttt{b\_class3\_phantom.py}. Thus
$|A_Q|=3^6$.
Every cocycle in $H^1(\Gamma,A_Q(-1))$ can again be
normalized at $c$ and is then determined by its single
value at $\tau$. Consequently this cohomology group
has at most $3^6$ elements. The torsion subgroup in
(B.21.2) has order $3^{59+2\cdot27+3\cdot7}=3^{134}$.
Its homomorphism to this finite target therefore has
a kernel of order at least $3^{128}$.

This is an obstruction to treating the fixed finite
quotient as complete information about $\eta$.
Additional geometric identities could constrain the
actual class to a much smaller subset. Neither the
three-dimensional path-comparison kernel of B.21
nor the finite Frobenius images alone supplies that
additional constraint.
\end{proof}

\subsection{A complete local-order law in class two}

\begin{statement}{B.23}{The finite-base arithmetic needed at three and five}{Proved}
Let $\ell=3$ or $5$, let $R=\mathbb Z_\ell$, and put
\[
 E:\ y^2=x^3+x,\qquad V=T_\ell E,\qquad
 S=\{2,\ell,\infty\},\qquad F=\mathbb Q(\zeta_8).
\]
Write $\chi_d$ for the quadratic character of $\mathbb Q(\sqrt d)$.
Let $U_{\rm CM}$ be the rank-two complement of
$R(1)\chi_{-1}$ in $\operatorname{Sym}^2V$, defined by the
two square CM characters. For each of
\[
 T\in\{V,\ V\chi_2,\ V(1),\ V(1)\chi_2,\ U_{\rm CM}\}
                                                               \tag{B.23.1}
\]
and every $n\geq1$ one has
\[
 H^0(G_{\mathbb Q,S},T/\ell^n)=
 H^2(G_{\mathbb Q,S},T/\ell^n)=0,\qquad
 H^1(G_{\mathbb Q,S},T/\ell^n)\simeq R/\ell^n.
                                                               \tag{B.23.2}
\]
The transition maps on $H^1$ are surjective. Consequently the
continuous groups are $H^1(G_{\mathbb Q,S},T)=R$ and
$H^2_{\rm cts}(G_{\mathbb Q,S},T)=0$.
Moreover $H^1(\mathbb Q,V)=H^1(G_{\mathbb Q,S},V)$.

For the twelve-cusp torus of the plus quartic, write
\[
 W=R(1)[\operatorname{Gal}(F/\mathbb Q)]^3/
             R(1)\,\mathbf 1_{\rm diag}.
\]
The actual rank-seventeen affine abelian coefficient is
$A=W\oplus V^3$. Its exterior square has the integral
decomposition
\[
 \begin{split}
 \Lambda^2A={}&B_W\oplus M\oplus B_T\oplus U_{\rm CM}^3,\\
 B_W={}&\Lambda^2W,\qquad
 M\simeq V(1)^{15}\oplus(V(1)\chi_2)^{18},\\
 B_T={}&R(1)^6\oplus(R(1)\chi_{-1})^3 .
 \end{split}                                                    \tag{B.23.3}
\]
For $B_{\rm rem}=B_W\oplus M\oplus U_{\rm CM}^3$,
$H^2(G_{\mathbb Q,S},B_{\rm rem}/\ell^n)=0$ for all $n$.
For the remaining summand the map
\[
 H^2_{\rm cts}(G_{\mathbb Q,S},B_T)
       \longrightarrow H^2_{\rm cts}(\mathbb Q,B_T)
                                                               \tag{B.23.4}
\]
is injective. These assertions concern the displayed actual
lattices, including their integral splittings.
\end{statement}
\begin{proof}
\emph{Residual division fields at three.}
Put $E_a:y^2=x^3+ax$, where $a=1,4$; the second is the
$\chi_2$-twist of the first. Let
\[
 P_a(t)=t^8+216a^3t^4-3888a^6,\qquad
 L_a=\mathbb Q(i,\theta),\quad P_a(\theta)=0.
\]
The point with
\[
 Y=\theta/3,\qquad
 X=\frac{12a^2\theta^2}{\theta^4+108a^3}
                                                               \tag{B.23.5}
\]
lies on $E_a$ and satisfies $[2](X,Y)=(X,-Y)$.
The point and its image $(-X,iY)$ under $[i]$ are independent
in $E_a[3]$, because $t^2+1$ is irreducible over $\mathbb F_3$.
Conversely their coordinates recover $\theta=3Y$ and $i$.
Thus $L_a$ is exactly the third-division field. For example,
the actual cyclotomic embedding is
\[
 \sqrt3=\frac{3X^2+3a}{2a},\qquad
 \zeta_3=\frac{-1+i\sqrt3}{2}.
                                                               \tag{B.23.6}
\]
The defining composita have degree sixteen. Unconditional
integer-ideal calculations give
\[
 \begin{array}{c|c|c|c|c}
 a&[L_a:\mathbb Q]&|\operatorname{disc}L_a|&
            h(L_a)&(e,f)\text{ at both }2,3\\ \hline
 1&16&2^{36}3^{14}&1&(8,2)\text{, one prime}\\
 4&16&2^{32}3^{14}&1&(8,2)\text{, one prime}.
 \end{array}                                                    \tag{B.23.7}
\]
Here and below these are reproducible certificates:
\texttt{b\_quartic\_fields.gp} constructs the displayed radical
fields, checks the point and division identities, and calls
\texttt{bnfcertify} with its full default certification.
All fifteen certificates return one; the output records every
polynomial, discriminant, class number, and local decomposition.
The meaning of full certification is documented in
\citet[\texttt{bnfcertify}]{BPARI2026}; no conditional class-group
bound is used.

The residual group is the full normalizer of the nonsplit
Cartan in $\operatorname{GL}_2(\mathbb F_3)$: its order is
sixteen, and on $\mathbb F_9$ it is generated by multiplication
by $1+i$ and conjugation. If $\tau$ is the first generator,
its two characters on $V$ have exponents $1,3$ modulo eight.
On $V(1)$ these are $5,7$, and on $U_{\rm CM}$ they are $2,6$.
Twisting by $(-1)$ adds four to these exponents. None is zero.
Both decomposition groups in (B.23.7) are the full group,
so the necessary global invariants of $T/\ell$ and local
invariants of $(T/\ell)(-1)$ vanish. The same calculation in
$L_4$ treats the quadratic twists.

\emph{Residual division fields at five.}
Let $K=\mathbb Q(i)$, and set $\sigma=-1$ for $a=1$ and
$\sigma=1$ for $a=4$. Choose
\[
 u^4=\sigma(1+2i),\qquad v^4=\sigma(1-2i),\qquad L=K(u,v).
                                                               \tag{B.23.8}
\]
The two valuations at $1+2i$ and $1-2i$ prove that these
Kummer classes are independent of order four. Hence
$[L:K]=16$. Complex conjugation exchanges $u,v$, so $L/\mathbb Q$
is Galois of degree thirty-two. Its $K$-group is
$C_4^2$, with $g_{r,s}(u,v)=(i^ru,i^sv)$.
For
\[
 Y=\begin{cases}(1+i)u/(2-i),&a=1,\\4u/(2-i),&a=4,\end{cases}
 \qquad X=\frac{5Y^2}{a(4+2i)},
                                                               \tag{B.23.9}
\]
direct substitution gives $Y^2=X^3+aX$ and
$[2](X,Y)=(-X,iY)=[i](X,Y)$.
This is a nonzero point killed by $[2-i]$; its conjugate is
killed by $[2+i]$. The two points span $E_a[5]$ and recover
$u,v$. Thus (B.23.8) is the actual fifth-division field.
On its two CM lines the characters are
$\alpha(g_{r,s})=2^r$ and $\beta(g_{r,s})=2^{-s}$, with
cyclotomic character $\alpha\beta$.

For $C_{j,k}=\ker((r,s)\mapsto jr+ks\bmod4)$, the fixed
fields $L^{C_{j,k}}=K(u^jv^k)$ have degree eight. The six
primitive characters up to multiplication by a unit are
\[
 (j,k)=(1,0),(0,1),(1,1),(1,2),(1,3),(2,1).
\]
The following entries are in that order:
\[
 \begin{array}{c|c|c}
 \sigma&\text{class numbers}&\text{absolute discriminants}\\ \hline
 -1&(1,1,2,2,1,2)&
 (32768000,32768000,64000000,\\
 &&\quad819200000,4000000,819200000)\\
 1&(1,1,2,2,1,2)&
 (8192000,8192000,64000000,\\
 &&\quad204800000,4000000,204800000).
 \end{array}                                                    \tag{B.23.10}
\]
These twelve certified fields have no five-primary class group.
Norm and extension on ideal classes, followed by averaging over
$C_{j,k}$, identify the $C_{j,k}$-fixed part of
$\operatorname{Cl}(L)\{5\}$ with the corresponding fixed-field
class group. Every character of $C_4^2$ is trivial on at least
one of these six kernels, including the characters of order two.
The idempotent decomposition over $\mathbb Z_5$ therefore gives
$\operatorname{Cl}(L)\{5\}=0$.

Only three of these quotients are needed for the local
invariants. At $2$ and $5$ their decompositions are
\[
 \begin{array}{c|c|c}
 (j,k)&2&5\\ \hline
 (1,0),(0,1)&(8,1)&(4,1),(1,4)\\
 (1,1)&(4,2)&(4,1),(4,1).
 \end{array}                                                    \tag{B.23.11}
\]
Every decomposition group in these cyclic quotients over $K$
has order four. The characters of $V(-1)$ are
$\alpha^{-1},\beta^{-1}$; those of $V(1)(-1)$ are
$\alpha,\beta$; those of $U_{\rm CM}(-1)$ are
$\alpha/\beta,\beta/\alpha$. Their kernels are precisely
among the three rows just tested. Thus none has local
invariants at $2$ or $5$. None of the untwisted modules
has global invariants. The calculation for $E_4$ supplies
the same assertions for the twists in (B.23.1).

\emph{From the finite fields to cohomology.}
Use the appropriate field $L$ above. It has degree prime
to $\ell$, contains $\mu_\ell$, trivializes $T/\ell$, and
has $\operatorname{Cl}(\mathcal O_L)\{\ell\}=0$.
The Kummer sequence over $\mathcal O_{L,S}$ identifies
$H^2(\mu_\ell)$ with the $\ell$-torsion of its Brauer group.
The local-invariant injection puts the latter inside the
permutation module on primes above $2,\ell$.
Tensoring by $(T/\ell)(-1)$ and taking the exact invariants
for $\operatorname{Gal}(L/\mathbb Q)$ gives zero, by the
local invariant calculations above. Restriction and averaging
therefore prove $H^2(G_{\mathbb Q,S},T/\ell)=0$.
The same argument gives $H^0=0$.
The sequences with successive quotients $T/\ell$ prove
both vanishing assertions for every $n$.
Complex conjugation has one positive and one negative
eigenline in each module in (B.23.1).
The global Euler characteristic
\citep[I, Theorem 5.1]{BMilne2006} gives
$|H^1(G_{\mathbb Q,S},T/\ell^n)|=\ell^n$.
The coefficient reduction maps on $H^1$ are onto because
the next $H^2$ is zero. The quotient by $\ell$ is
one-dimensional, so these groups are cyclic of the indicated
orders, with compatible generators.
Finiteness kills $\varprojlim^1 H^1$, giving (B.23.2) in
continuous cohomology. At a good prime $p\ne\ell$,
$V(-1)^{F_p}=0$, by the two elliptic Frobenius weights.
The inertia of any integral $V$-cocycle is therefore zero;
this proves the asserted equality of full and unramified $H^1$.

\emph{The integral quartic summands.}
The three morphisms from $x^4+y^4+z^4=0$ to $E$ are
\[
 \begin{split}
 f_0[x:y:z]&=[-x^2y:z^2x:y^3],\\
 f_1[x:y:z]&=[-y^2z:x^2y:z^3],\\
 f_2[x:y:z]&=[-z^2x:y^2z:x^3].
 \end{split}                                                    \tag{B.23.12}
\]
The coordinates have no common zero on the curve and satisfy
the homogeneous Weierstrass equation. Each map has degree
four; its differential is twice the corresponding basis
differential. Their distinct sign characters give
$f_{i*}f_j^*=0$ for $i\ne j$, while $f_{i*}f_i^*=[4]$.
Thus $\Phi=(f_{0*},f_{1*},f_{2*})$ and
$\Psi=\sum_i f_i^*$ satisfy both $\Phi\Psi=[4]$ and
$\Psi\Phi=[4]$. The latter follows on differentials, hence
on homomorphisms in characteristic zero.
Since four is a unit in $R$, this is an integral
identification of the compact Tate module with $V^3$.
The tangent divisors are $4d-4O$, so pullback by $[4]$
splits the generalized-Jacobian extension. Its unique
map $j$ with $pj=[4]$, divided by four, gives the actual
equivariant splitting $A=W\oplus V^3$.
All polynomial and affine-chart checks for (B.23.12) are
recomputed in \texttt{b\_local\_order\_law.py}.

The four character idempotents of $\operatorname{Gal}(F/\mathbb Q)$
give
\[
 W=R(1)^2\oplus(R(1)\chi_{-1})^3
       \oplus(R(1)\chi_2)^3\oplus(R(1)\chi_{-2})^3.
\]
The automorphism $I=[i]$ identifies $V\chi_{-1}$ with $V$.
Decomposing $\Lambda^2(V^3)$ using division by two,
and projecting $\operatorname{Sym}^2V$ with
$I\otimes I$, gives (B.23.3) over $R$ itself.
Thus all its projections and sections have denominators
only powers of two. In particular no division by three
enters the coefficient decomposition when $\ell=3$.

It remains to treat $B_W$. Over $F$ it is $R(2)^{55}$.
The exact result used here is
$K_2(\mathcal O_F)=0$, the tame-kernel theorem of
\citet[Theorem 3.5]{BGuoQinSong2003}.
It is not a statement that $K_2(F)$ vanishes.
Localization to $\mathcal O_F[1/(2\ell)]$ introduces only
residue multiplicative groups at $2,\ell$, of order prime
to $\ell$. Consequently its $K_2/\ell^n$ is zero.
The integral comparison
\citep[Proposition A2.1]{BKahn1993}
identifies this quotient with $H^2(\mu_{\ell^n}^{\otimes2})$.
Restriction and corestriction of degree four prove the
claimed vanishing for $B_W$, and (B.23.2) handles the
other summands of $B_{\rm rem}$.
Finally $F$ has class number one (also certified in the
script). Kummer theory makes the finite-level $H^2$ of
$\mu_{\ell^n}$ over its $S$-integer ring inject into
the field Brauer group. Its unit $H^1$ is finite at
every exponent, and the full field Kummer $H^1$
transitions are surjective. Passing to the continuous
limit preserves this injection. Restriction and
corestriction give (B.23.4).
\end{proof}

\begin{statement}{B.24}{The unified integral class-two local-order law}{Proved}
For the cubic at $\ell=3$ and the plus quartic at
$\ell\in\{3,5\}$, the image of the actual class-two section
set in its abelian input space is the subset described below.
In every case the rule has the same form: every displayed
local correction is integral, and its minimum coordinate
valuation tends to $+\infty$ outside finite sets of primes.
Valuation of zero means $+\infty$.

\emph{Cubic input and its forty-four scalar corrections.}
Use $O=C_0$, $K=\mathbb Q(\zeta_3)$, and the ordered basis
\[
 (f_1,f_2,A_0,A_1,A_2,B_0,B_1,B_2,C_1,C_2)
                                                               \tag{B.24.1}
\]
of the actual abelian lattice $V=T_3J_{E,D}$. Here the last
eight symbols denote positive peripheral cocharacters, and
\[
 A_i=[0:-\zeta_3^i:1],\quad
 B_i=[-\zeta_3^i:0:1],\quad C_i=[1:-\zeta_3^i:0].
\]
Write $W$ for their span and let $U_1,U_2\in W$ be the
columns of
\[
 U=\begin{pmatrix}
 0&-1\\2&-2\\1&0\\0&-2\\1&-1\\2&0\\1&-2\\2&-1
 \end{pmatrix},\qquad
 R_3=\begin{pmatrix}0&-1\\1&-1\end{pmatrix}.
                                                               \tag{B.24.2}
\]
The compact inclusion has block column $\binom{3I_2}{U}$.
Put $w=(3f_1+U_1)\wedge(3f_2+U_2)$, whose coefficient at
$A_0\wedge A_2$ is one.

Let $\alpha\in H^1(\mathbb Q,V)$. For a prime
$\mathfrak q\nmid3$ of $K$, put $q=N\mathfrak q$.
Its elliptic Frobenius is
$A_E=aI_2+bR_3$, where $a+b\zeta_3$ is the primary
Jacobi-sum generator, congruent to one modulo three.
At an inert rational prime $p$ it is $-p$ and $q=p^2$.
The actual affine matrix is
\[
 A_{\mathfrak q}=
 \begin{pmatrix}
 A_E&0\\ U(A_E-qI_2)/3&qI_8
 \end{pmatrix}.                                                \tag{B.24.3}
\]
Take a tame generator $t_{\mathfrak q}$ and arithmetic
Frobenius $F_{\mathfrak q}$, with $FtF^{-1}=t^q$.
Set $a_{\mathfrak q}=\alpha(F_{\mathfrak q})$ and
$v_{\mathfrak q}=\alpha(t_{\mathfrak q})
 =\sum_d b_{\mathfrak q,d}e_d\in W$.
Let $k_{\mathfrak q,d}$ be an integral representative of
the unramified Kummer class of the following actual framed
point of $J_{E,D}$: $z_d=\mathcal O(d-O)$ with frames
$s_d^{-1}$ at $d$, $s_O$ at $O$, and one elsewhere.
The parameters are $s_{A_i}=x/z$, $s_{B_i}=y/z$ and
$s_{C_i}=z/x$. All the finite-field input is given by
\[
 \begin{split}
 3z_d&=\left(
   \operatorname{lead}_r(L_d/L_O)/
   \operatorname{lead}_O(L_d/L_O)\right)_{r\ne O},\\
 L_{A_i}&=y+\zeta_3^iz,\quad
 L_{B_i}=x+\zeta_3^iz,\quad
 L_{C_i}=x+\zeta_3^{2i}y,\quad L_O=x+y.
 \end{split}                                                    \tag{B.24.4}
\]
At $d$ the leading value divides by $s_d^3$, and at
$O$ it multiplies by $s_O^3$; elsewhere it is evaluation.
The Kummer class is read from divisions of $\overline z_d$
in $J_{E,D}(\overline{\mathbb F}_q)$ by applying $F-1$.
It lies in the finite quotient $V/(A_{\mathfrak q}-1)V$.

Form, in the basis (B.24.1),
\[
 \begin{split}
 D_{\mathfrak q}
 &=a_{\mathfrak q}\wedge v_{\mathfrak q}
       -\sum_d b_{\mathfrak q,d} k_{\mathfrak q,d}\wedge e_d\\
 &=f_1\wedge X_1+f_2\wedge X_2+Y,
       \qquad X_i\in W,\quad Y\in\Lambda^2W .
 \end{split}                                                    \tag{B.24.5}
\]
Regard $X=(X_1,X_2)^{\rm t}$ as a two-by-eight matrix.
The following operations take place over $\mathbb Q_3$:
\[
 \begin{split}
 Z&=(A_E-I_2)^{-1}X,\qquad Z=(Z_1,Z_2)^{\rm t},\\
 Y_*&=\frac{Y-\tfrac13(U_1\wedge X_1+U_2\wedge X_2)}{q-1}
          +\tfrac13(U_1\wedge Z_1+U_2\wedge Z_2),\\
 z_*&=(Y_*)_{A_0\wedge A_2}.
 \end{split}                                                    \tag{B.24.6}
\]
Define $\rho_{\mathfrak q}(\alpha)$ to be the minimum
$3$-adic valuation of these forty-four coordinates:
\[
 \begin{array}{c|c}
 \text{coordinate block}&\text{correction}\\ \hline
 f_1\wedge f_2&9z_*\\
 f_1\wedge W&-Z_1+3z_*U_2\\
 f_2\wedge W&-Z_2-3z_*U_1\\
 \Lambda^2W\text{ except }A_0\wedge A_2&
                   -Y_*+z_*U_1\wedge U_2 .
 \end{array}                                                    \tag{B.24.7}
\]
The complete cubic image is
\[
 \left\{\alpha:\ 
   \rho_{\mathfrak q}(\alpha)\geq0\text{ for all }\mathfrak q\nmid3,
   \quad \rho_{\mathfrak q}(\alpha)\longrightarrow+\infty\right\}.
                                                               \tag{B.24.8}
\]
This is a two-by-two formula for the full intrinsic
class-two quotient, with no omission of the integral
index-nine compact gluing.

\emph{Quartic input and its scalar order functions.}
Use the plus quartic $x^4+y^4+z^4=0$ and
$F=\mathbb Q(\zeta_8)$. Let $\alpha_i=\zeta_8^{1+2i}$,
and take cusps
$A_i=[0:\alpha_i:1]$, $B_i=[\alpha_i:0:1]$,
$C_i=[\alpha_i:1:0]$, with $O=B_0$.
Use unit parameters $x/z,y/z,z/y$, respectively.
The abelian section origin is the average of the four
conjugate $B_i$ tangents. It has an actual class-two lift.
In the splitting of B.23, every input has unique coordinates
\[
 (b,\mathbf r)\in H^1(\mathbb Q,W)\oplus R^3,\qquad
 e_{\mathbf r}=\sum_{i=1}^3r_i j_i h,
                                                               \tag{B.24.9}
\]
where $h$ is any fixed generator of $H^1(\mathbb Q,V)=R$
and $j_i:V\to A$ are the actual compact inclusions.
Equivalently the first factor is
\[
 H^1(\mathbb Q,W)=
 (\widehat{F^\times}_{\,\ell})^2
 \times\ker\left(N_{F/\mathbb Q}:
      \widehat{F^\times}_{\,\ell}\to
      \widehat{\mathbb Q^\times}_{\,\ell}\right).
                                                               \tag{B.24.10}
\]

Restrict $b$ to $F$ in the order
$(A_0,A_1,A_2,A_3,B_1,B_2,B_3,C_0,C_1,C_2,C_3)$.
Let $\mathbf 1$ denote the length-four all-one column and put
\[
 \begin{gathered}
 H=\begin{pmatrix}
 -8&-3&-2&-3\\-3&-8&-3&-2\\
 -2&-3&-8&-3\\-3&-2&-3&-8
 \end{pmatrix},\quad r=\begin{pmatrix}-5\\-6\\-5\end{pmatrix},\\
 B=\begin{pmatrix}-10&-6&-4\\-6&-12&-6\\-4&-6&-10\end{pmatrix},
 \qquad
 N=\begin{pmatrix}
 H&{\bf1}r^{\rm t}&-4{\bf1}{\bf1}^{\rm t}\\
 r{\bf1}^{\rm t}&B&r{\bf1}^{\rm t}\\
 -4{\bf1}{\bf1}^{\rm t}&{\bf1}r^{\rm t}&H
 \end{pmatrix},\\
 v=(-\tfrac12,-\tfrac12,-\tfrac12,-\tfrac12,
       -\tfrac58,-\tfrac34,-\tfrac58,
       -\tfrac12,-\tfrac12,-\tfrac12,-\tfrac12)^{\rm t}.
 \end{gathered}                                                 \tag{B.24.11}
\]
For every $i<j$ let $k=\kappa_\ell(2)$ and define
\[
 p_{ij}=v_i-N_{ij}/8,\quad q_{ij}=v_j-N_{ij}/8,\quad
 X_{ij}=b_i+p_{ij}k,\quad Y_{ij}=b_j+q_{ij}k .
                                                               \tag{B.24.12}
\]
This gives fifty-five pairs of explicit Kummer classes.

Here is their scalar test on all completed inputs.
Take $\epsilon=1+\sqrt2$ and a generator $\pi_{\mathfrak p}$
of every prime ideal of $F$. Every Kummer class has unique
coordinates
\[
 x=a_\epsilon\kappa_\ell(\epsilon)
       +\sum_{\mathfrak p}a_{\mathfrak p}
                         \kappa_\ell(\pi_{\mathfrak p}),
 \qquad a_{\mathfrak p}\longrightarrow0.
                                                               \tag{B.24.13}
\]
For a pair $X,Y$ denote its two coefficient families by
$a,b$. At $\mathfrak q\nmid\ell$, put
$Q=N\mathfrak q$, $e=\nu_\ell(Q-1)$, and choose
a primitive $\ell^e$-th root $\eta_{\mathfrak q}$
in its residue field. For $u=\epsilon$ or
$u=\pi_{\mathfrak p}$ with $\mathfrak p\ne\mathfrak q$,
let $r_{\mathfrak q,u}$ be the integer in $[0,\ell^e)$
determined by
\[
 \overline u^{(Q-1)/\ell^e}
       =\eta_{\mathfrak q}^{\,r_{\mathfrak q,u}}.
\]
Set
\[
 \begin{split}
 A_{\mathfrak q}
 &=a_\epsilon r_{\mathfrak q,\epsilon}
    +\sum_{\mathfrak p\ne\mathfrak q}
                    a_{\mathfrak p}r_{\mathfrak q,\pi_{\mathfrak p}},\\
 B_{\mathfrak q}
 &=b_\epsilon r_{\mathfrak q,\epsilon}
    +\sum_{\mathfrak p\ne\mathfrak q}
                    b_{\mathfrak p}r_{\mathfrak q,\pi_{\mathfrak p}},\\
 \rho^{ij}_{\mathfrak q}(b)
 &=\nu_\ell(A_{\mathfrak q}b_{\mathfrak q}
                -B_{\mathfrak q}a_{\mathfrak q})-e .
 \end{split}                                                    \tag{B.24.14}
\]
If $e=0$ take all $r_{\mathfrak q,u}=0$ and
$\rho^{ij}_{\mathfrak q}=+\infty$.
The toric test is
\[
 \rho^{ij}_{\mathfrak q}(b)\geq0\text{ for every }\mathfrak q\nmid\ell,
 \qquad
 \rho^{ij}_{\mathfrak q}(b)\longrightarrow+\infty
          \quad\text{for all }i<j.
                                                               \tag{B.24.15}
\]

There is one further order function for the mixed part.
At a rational $p\notin S$, put $v_p=b(t_p)\in W$ and
\[
 \overline u_p=
  (F_{E,p}-I)^{-1}h(F_p)\pmod V,\qquad
 \operatorname{ord}(\overline u_p)=\ell^{m_p}.
                                                               \tag{B.24.16}
\]
The class is in $(V\otimes\mathbb Q_\ell/V)^{F_p}$,
identified with $E(\mathbb F_p)[\ell^\infty]$. In particular
$m_p\leq\nu_\ell(\#E(\mathbb F_p))$. For $\mathbf r\ne0$
write $t=\min_i\nu_\ell(r_i)$ and set
\[
 \rho^{\rm mix}_p(b,\mathbf r)
  =\nu_\ell(v_p)+t-m_p\quad(m_p>0),\qquad
 \rho^{\rm mix}_p=+\infty\quad(m_p=0).
                                                               \tag{B.24.17}
\]
For $\mathbf r=0$ set every $\rho^{\rm mix}_p=+\infty$.
The mixed test is
\[
 \rho^{\rm mix}_p(b,\mathbf r)\geq0\text{ for all }p\notin S,
 \qquad
 \rho^{\rm mix}_p(b,\mathbf r)\longrightarrow+\infty .
                                                               \tag{B.24.18}
\]
The complete quartic image consists exactly of the pairs
$(b,\mathbf r)$ satisfying both (B.24.15) and (B.24.18).
Thus the function cutting out the image is fully specified
by scalar residue logarithms and the orders of the actual
local Kummer classes (B.24.16). Replacing $m_p$ by the
possibly larger valuation of the group order would change
the criterion.

All tests are independent of cochain representatives,
compatible path choices, and Frobenius prolongations.
The finite logarithms in (B.24.14) are used modulo their
specified powers; they are not asserted to equal full
$\ell$-adic Frobenius values. Every nonempty fibre is a
torsor under the actual central group
$H^1_{\rm cts}(\mathbb Q,\Lambda^2 A)$, with $A=V$
in the cubic case.
\end{statement}
\begin{proof}
\emph{The continuous localization mechanism.}
For a finite free unramified coefficient $T$ on a fixed
$S$-integer base, suppose
$H^2(G_S,T/\ell^n)=H^3(G_S,T/\ell^n)=0$ and the finite-base
$H^1$ is finite at every $n$. First take the direct limit
of finite-coefficient localization over finite sets of
removed primes. The support complex has two terms
\[
 \bigoplus_{\mathfrak q\notin S} T/\ell^n
   \ \xrightarrow{\ q^{-1}F_{T,\mathfrak q}-1\ }\
 \bigoplus_{\mathfrak q\notin S} T/\ell^n .
\]
Both systems of terms have surjective reduction maps.
Their derived inverse limit is therefore the literal
two-term complex on
$\mathcal R(T)=\{(z_{\mathfrak q}):z_{\mathfrak q}\in T,\
z_{\mathfrak q}\to0\}$. The finite-base vanishings and
finiteness kill the adjacent derived limits. Thus
\[
 H^2_{\rm cts}(G,T)=
 \mathcal R(T)/
 (q^{-1}F_{T,\mathfrak q}-1)_{\mathfrak q}\mathcal R(T).
                                                               \tag{B.24.19}
\]
The image is its literal image, not its closure.
These are the localization and inverse-system comparison
maps of \citet[II.0.1, II.1.7, II.2.9]{BMilne2006}
and \citet[Proposition 1.6, Theorem 2.2]{BJannsen1988}.
This proves the common origin of integrality and decay
in the two cases.

\emph{Cubic geometry and the inverse formula.}
The nine cusps are $E[3]$. With $A_0,A_2$ as the ordered
three-torsion basis their Weil pairing is $\zeta_3$.
Their coordinates in this basis are
$(1,0),(2,2),(0,1),(2,0),(1,1),(0,2),(2,1),(1,2)$.
Pulling back the generalized-Jacobian extension by $[3]$
trivializes its characters $\mathcal O(d-O)$.
The vector frame transforms by the inverse Weil pairing,
so its gluing row is $(v_d,-u_d)$, giving (B.24.2).
Equivariance of the columns $\binom{3I_2}{U}$ gives
(B.24.3). The normalized tangent-line divisors give
(B.24.4). At the two exceptional evaluations the leading
values are, respectively,
\[
 \begin{array}{c|cc}
 d&\operatorname{lead}_d(L_d/L_O)&
           \operatorname{lead}_O(L_d/L_O)\\ \hline
 A_i&1/3&3\\
 B_i&1/3&-3\\
 C_i,\ i=1,2&-1/(3(1-\zeta_3^i))&
                       -3(1-\zeta_3^{2i}).
 \end{array}
\]
They are units away from three. Passing from the endpoint
frame of a tangent path to this divisor frame introduces
only signs of order two, whose pro-three Kummer torsor is
canonically trivial. Thus $k_d$ is the actual abelian
path cocycle, with its stated sign.

The good smooth pair and its unit tangent give unramified
based pro-three action outside three. The relevant
prime-to-residue-characteristic specialization, including
inertia identification, is
\citet[Expos\'e XIII, \S\S2.9--2.12]{BSGA1}.
The tame relation then forces the inertia of $\alpha$
into $W$: its elliptic projection would be an eigenvector
of $F_E$ with eigenvalue $q$, which does not exist.
In class-two BCH coordinates the group law is
$(x,z)(y,t)=(x+y,z+t+x\wedge y/2)$.
On a peripheral generator the actual arithmetic linear
term is
$D_g(e_d)=-\chi(g)k_d(g)\wedge e_d$.
Compute $FtF^{-1}=t^q$ in the pullback extension over
$\alpha$. The central Frobenius coordinate cancels and
the equation for its inertia coordinate $m$ is
\[
 (q^{-1}\Lambda^2 A_{\mathfrak q}-1)m=-D_{\mathfrak q}.
                                                               \tag{B.24.20}
\]
This identifies the actual support representative, not
just its image in a finite quotient.

The kernel of the operator in (B.24.20) is the primitive
line $\mathbb Z_3w$. All its other eigenvalues are
nonzero by their Frobenius weights. Set
$g_i=f_i+U_i/3$. In this rational basis the operator is
$A_E-I_2$ on $g_i\wedge W$ and $q-1$ on $\Lambda^2W$.
Writing (B.24.5) in that basis and inverting gives
$m_0=\sum_i f_i\wedge Z_i+Y_*$ for the positive
right-hand side $D_{\mathfrak q}$.
Subtracting $z_*w$ eliminates the pivot $A_0\wedge A_2$.
Negating this representative gives exactly (B.24.7).
Thus these are the integral coordinates of the unique
solution in $\Lambda^2V/\mathbb Z_3w$, not those of an
unjustified direct sum with the compact line.
The determinant of the resulting forty-four-dimensional
map, up to a unit, is
$\#E(\mathbb F_q)^8(q-1)^{28}/9$.
The factor nine is retained by the $9z_*$ term.
The script independently verifies both inverse identities
on all forty-four basis columns for six distinct elliptic
Frobenius matrices, in exact rational arithmetic.

For clarity the finite-base cubic vanishing has a direct
Kummer proof. The ring $\mathbb Z[\zeta_3]$ is Euclidean,
and there is exactly one prime above three.
Thus $\operatorname{Pic}(\mathcal O_K[1/3])=0$ and
its Brauer group is zero: a Brauer class can have
only one finite invariant, whose sum must be zero.
As $\mu_3\subset K$, Kummer theory gives
$H^2(G_{K,\{3\}},\mathbb F_3)=0$.
The matrices (B.24.3) are block unipotent modulo three,
so the residual action has three-group image.
A composition series therefore
gives the vanishings needed in (B.24.19) for
$T=\Lambda^2V$ and every exponent.
Continuity gives $v_{\mathfrak q}\to0$, and all other
entries of (B.24.5) are integral and bounded. Hence its
support family is restricted. The fixed integral
pivot section of $\Lambda^2V/\mathbb Z_3w$ lifts
restricted families to restricted families.
Equation (B.24.19) now proves both directions of
(B.24.8). Restriction from $\mathbb Q$ to $K$ is
injective on the obstruction group, since its
corestriction composite is multiplication by two.

\emph{Quartic tangent origin and all peripheral terms.}
The tangent functions are
$L_{A_i}=y-\alpha_i z$, $L_{B_i}=x-\alpha_i z$,
$L_{C_i}=x-\alpha_i y$.
The ratios $L_d/L_O$ have divisors $4d-4O$. At a cusp $d$, the leading
coefficient of $L_d$ in its specified parameter is
$\alpha_i/4$. Evaluating all ratios of leading terms
of $L_d/L_O$ gives
\[
 4z_d=(\zeta_8^{K_{rd}}(\sqrt2)^{N_{rd}})_{r\ne O}.
                                                               \tag{B.24.21}
\]
The independently recomputed matrix $N$ is exactly
(B.24.11); the script checks all 121 equalities in
$\mathbb Q[t]/(t^4+1)$. Roots of unity have zero
odd-primary Kummer class, so
$\kappa(z_d)=N_{\bullet d}\kappa_\ell(2)/8$.
This fixes every path contribution integrally.
If $L(e_d)=e_d\wedge N_{\bullet d}/8$, the actual
restricted action on $W\oplus\Lambda^2W$ is
\[
 g(x,z)=\bigl(\chi(g)x,\
        \chi(g)^2z+\chi(g)k(g)L(x)\bigr).
                                                               \tag{B.24.22}
\]
The average of the four conjugate $B_i$ tangent classes
has abelian coordinate $k v$; summing the three nonzero
columns in (B.24.21) and dividing by four gives
the stated $v$. Its central coordinate
$k(g)^2L(v)/2$ satisfies the class-two cocycle identity
in (B.24.22). This proves the existence of its
class-two lift over $F$.
Here is the exact averaging construction used to descend.
In an extension with abelian kernel $T$, choose four
coset representatives $g_i$ for $G/H$, lifts $e_i$,
and a splitting $s_H$ over $H$.
For an element $e$ above $g$, write
$gg_i=g_{gi}h_i(g)$ and set
\[
 B(e)=\sum_i
       e\,e_i\,s_H(h_i(g))^{-1}e_{gi}^{-1}\ \in T .
\]
Multiplying the factors and reindexing gives
$B(ef)=B(e)+gB(f)$ and $B(z)=4z$ for $z\in T$.
Hence $(-B(e)/4)e$ is an independent-of-lift section.
Apply this first to the abelian extension; its restricted
class is the displayed average $kv$.
Apply it a second time to the pulled-back central
extension with its just-constructed $H$-splitting.
This gives an actual class-two lift over $\mathbb Q$
preserving that abelian origin. No nonabelian averaging
has been used.
Both constructions are defined over $G_{\mathbb Q,S}$.
Indeed the pair has good reduction outside two and
the tangent leading constants are units there;
the same specialization result as above applies at
every $p\notin S$. Locally the image of arithmetic
inertia in the coefficientwise tangential splitting
centralizes the full free pro-$\ell$ geometric group $N$.
Different such splittings have the same inertia image:
their ratio lies in $Z(N)=1$.
The closed subgroup generated by these inertia images
is therefore normal in the full arithmetic extension,
centralizes $N$, and intersects $N$ trivially.
Its arithmetic image is exactly the kernel of
$G_{\mathbb Q}\to G_{\mathbb Q,S}$.
Quotienting by it gives the required extension over
$G_{\mathbb Q,S}$ with unchanged geometric kernel.

In the shifted origin, direct substitution into the BCH
cocycle identity gives for the $(i,j)$ peripheral
obstruction
$X_{ij}\cup Y_{ij}+d\Phi_{ij}$, where
\[
 \Phi_{ij}(g)=
 \tfrac12 b_i(g)b_j(g)+q_{ij}k(g)b_i(g)
                 +\tfrac12 p_{ij}q_{ij}k(g)^2.
                                                               \tag{B.24.23}
\]
This equality is checked as a polynomial identity in
the two cocycle values and $\chi(g)$ in the script.
The coefficient splitting and Shapiro's cocycle maps
give (B.24.10). Class number one and
$\mathcal O_F^\times=\mu_8\langle1+\sqrt2\rangle$
give the coordinates (B.24.13); the unit assertion
is also checked using the certified fundamental
units of the same quartic field.

For $R(2)$ the localization differential is $Q-1$.
The tame relation for $X\cup Y$ gives its support
representative
$A_{\mathfrak q}b_{\mathfrak q}
 -B_{\mathfrak q}a_{\mathfrak q}$, with an immaterial
common sign. The finite residue logarithms determine
each unramified Kummer value modulo $\ell^e$.
Replacing the full values by these representatives
changes its correction by an integral family tending
to zero: the differences are in
$(Q-1)\mathbb Z_\ell$, and they are multiplied by
$a_{\mathfrak q},b_{\mathfrak q}\to0$.
All infinite sums converge because their coefficients
tend to zero and the integer residues are integral.
B.23 gives the finite-base $R(2)$ vanishing. Applying
(B.24.19) therefore proves precisely (B.24.15), with
its decay condition. Restriction/corestriction of
degree four returns the answer to $\mathbb Q$.

\emph{Why no unevaluated compact action remains in the law.}
Write the actual linear class-two action as
$g(x,z)=(\rho_A(g)x,\rho_B(g)z+D_gx)$ and
$c_g=D_g\rho_A(g)^{-1}$, a cocycle in
$\operatorname{Hom}(A,\Lambda^2A)$.
For each Tate projection
$\lambda:\Lambda^2A\to L$, $L=R(1)$ or
$R(1)\chi_{-1}$, the perfect pairings
\[
 B_\lambda(x,y)=\langle x,y\rangle
       \quad\text{or}\quad \langle Ix,y\rangle
                                                               \tag{B.24.24}
\]
identify $\operatorname{Hom}(V,L)$ with $V$.
The first is alternating; the second is symmetric,
because $I$ is symplectic and $I^2=-1$.
Consequently $[\lambda c j_i]$ is a scalar multiple
of $[h]$, by B.23. Its evaluation on $g h(t)$ is a
scalar multiple of the corresponding self-cup, modulo
an explicit coboundary.

Both relevant self-cups vanish. The symmetric one
has primitive $-B_\lambda(h,h)/2$.
For the alternating one, at $2$ every point of
$E(\mathbb Q_2)$ reduces into the nonsingular
subgroup of order two: an integral odd $x$ would make
$x(x^2+1)$ have odd valuation one, and a nonintegral
point reduces to infinity. The formal group kernel
is pro-two, so $E(\mathbb Q_2)[\ell^n]=0$.
Local duality and the prime-to-two Euler characteristic
\citep[I, Theorem 2.1 and Lemma 2.9]{BMilne2006}
give $H^1(\mathbb Q_2,E[\ell^n])=0$.
At good primes other than $\ell$, $h$ is unramified,
so its finite cup factors through a procyclic group
of $\ell$-cohomological dimension one. The real
contribution is zero. The Brauer reciprocity sequence
then kills the only remaining contribution, at $\ell$
\citep[VIII, Theorem 4.2]{BMilneCFT2020}.
The full-field Kummer $H^1$ transitions are onto,
so the continuous alternating cup also vanishes.
The injection (B.23.4) gives its vanishing on $G_S$.

All other coefficient projections have zero finite-base
$H^2$ by B.23. Thus every pure compact input
$e_{\mathbf r}$ has an actual continuous class-two
lift over $G_S$, including its linear arithmetic
term $D_g e_{\mathbf r}(t)$.
This argument neither assigns a value to the Ceresa
class nor assumes it to be zero.
For a pure toric input, restriction to $F$ puts the
entire obstruction in $B_W$ by (B.24.22), so
its other projections are boundaries by degree-four
restriction/corestriction.
Finally the full obstruction for $b+e_{\mathbf r}$
differs from the sum of the two pure obstructions
and the ordered mixed cup
$b(g)\wedge g e_{\mathbf r}(t)$ by
$d(b\wedge e_{\mathbf r})/2$.
After subtracting the just-constructed primitives,
the only remaining terms lie in the two distinct
summands $B_W$ and $M$.

\emph{The mixed local order and its sufficiency.}
Apply (B.24.19) to $M$, using B.23.
The tame relation in the ordered-cup extension has
central defect
$-p\,v_p\wedge e_{\mathbf r}(F_p)$.
Dividing by $p$ in the support map, its correcting
equation is
\[
 (p^{-1}F_{M,p}-1)c_p
        =v_p\wedge e_{\mathbf r}(F_p).
                                                               \tag{B.24.25}
\]
Write $F_{W,p}=pP_p$. The inertia cocycle relation
gives $P_pv_p=v_p$, so the unique rational solution is
\[
 c_p=v_p\wedge
          j_{\mathbf r}(F_{E,p}-I)^{-1}h(F_p).
                                                               \tag{B.24.26}
\]
Uniqueness follows also from the Frobenius weights:
every eigenvalue of $p^{-1}F_M$ has absolute value
$\sqrt p$, so none is one.
The boundary map in
$0\to V\to V\otimes\mathbb Q_\ell
       \to V\otimes\mathbb Q_\ell/V\to0$
identifies (B.24.16) with the unramified local
Kummer class of $h$. Choose a fractional representative
of $\overline u_p$ with coordinates in
$\mathbb Z[1/\ell]\cap[0,1)$.
It has minimum valuation $-m_p$ if $m_p>0$,
and is zero otherwise. The change from the actual
inverse in (B.24.26) is integral; multiplied by
$v_p\to0$, it is an integral restricted family.
In the integral tensor lattice
$M=W\otimes R^3\otimes V$ the valuation of a nonzero
pure tensor is the sum of its factor valuations.
This gives exactly (B.24.17)--(B.24.18).

For explicit finite evaluation of $m_p$, it suffices
to compute the marked cocycle $h$ and elliptic
Frobenius modulo $\ell^a$ with
$a=\nu_\ell(\#E(\mathbb F_p))$.
No uncomputed infinite precision is required.
Compatible representatives $h_n$ can, if needed,
be constructed by enumerating finite Galois
extensions unramified outside $S$ containing
$\mathbb Q(E[\ell^n])$, solving their finite cocycle
equations, and choosing a primitive class lifting
$h_{n-1}$. Such an enumeration is effective by
defining polynomials and discriminant tests.
B.23 proves that the required primitive class
exists and that the coefficient maps are onto,
so each finite stage terminates.
This prescription retains the actual representation
and Frobenius marking; it does not replace $h$
by an arbitrary homomorphism.

If all tests hold, the correcting families belong
to $\mathcal R$ in (B.24.19), which provides an
actual continuous primitive. One may also construct
it exponent by exponent: the prescribed corrections
have finite support modulo each $\ell^n$;
localization supplies a primitive with those
inertia values. The remaining choices form a finite
torsor under unramified one-cocycles.
Every finite collection of compatibility conditions
can be met at its largest exponent, so their
inverse limit supplies a continuous primitive.
The boundary corrections above return an actual
class-two section. Conversely any section gives
these same corrections by its tame relation.
This proves necessity as well as sufficiency.
Changes of representatives alter the corrections
by integral restricted families, as the displayed
local equations show. Finally $H^0(\mathbb Q,A)=0$
by Frobenius weights. A difference of two aligned
central section coordinates is a one-cocycle,
and central conjugacy is its ordinary coboundary.
This proves the fibre assertion.
\end{proof}

\begin{statement}{B.25}{Unbounded local orders and explicit failure of additivity}{Proved}
The class-two images in B.24 are not subgroups of their
displayed abelian input spaces.
For the cubic, set $P=A_0$, $Q=B_0$, and let $\alpha_Q$
be the class of the unit tangent at $Q$ relative to the
unit tangent at $O=C_0$. The two inputs
\[
 a=\kappa_3(3)e_P,\qquad
 b=\alpha_Q+\kappa_3(13)e_Q
                                                               \tag{B.25.1}
\]
both lift to the actual intrinsic class-two extension;
$a+b$ does not lift, already locally at $13$.

For the quartic at either $\ell=3$ or $5$, the local
orders $m_p$ in (B.24.16) are unbounded.
More precisely, for every $n\geq1$ there are infinitely
many rational primes $p$ such that
\[
 \begin{gathered}
 p\equiv1\pmod8,\qquad \nu_\ell(p-1)=n,\qquad
 2\in(\mathbb F_p^\times)^{\ell^n},\\
 F_{E,p}\equiv I\pmod{\ell^n},\qquad
 h(F_p)\notin\ell V,\qquad m_p\geq n .
 \end{gathered}                                                 \tag{B.25.2}
\]
For any nonzero compact vector $\mathbf r$, choose such a
prime with $n>\min_i\nu_\ell(r_i)$.
Then the pure compact input $(0,\mathbf r)$ and the
pure toric input
\[
 \left(\kappa_\ell(p)\sum_{i=0}^3e_{A_i},\,0\right)
                                                               \tag{B.25.3}
\]
both lift, whereas their sum fails the mixed divisibility
test at $p$. Equivalently (B.25.3) has input
$(p,1,1)$ in the three Kummer factors of (B.24.10).
\end{statement}
\begin{proof}
The first cubic input is unramified outside three;
the finite-base $H^2$ vanishing in B.24 gives its lift.
The second is a section of the actual cusp decomposition
group at $Q$, so it lifts even to the full geometric
pro-three quotient. In the local equation (B.24.5) for
their sum, choose the path defining $\alpha_Q$ to compute
$k_Q$. The path terms cancel, leaving
$\kappa_3(3)(F_{13})\,e_P\wedge e_Q$.
At either prime of $K$ above $13$, $q-1=12$ has valuation
one and $\kappa_3(3)(F_{13})$ is a unit: $3$ is not a
cube modulo $13$, since $3^4\equiv3\ne1$.
The quotient correction has the nonintegral coordinate
$-\kappa_3(3)(F_{13})/12$ at $e_P\wedge e_Q$.
This wedge is distinct from the deleted
$A_0\wedge A_2$ pivot, so no compact-line adjustment
can remove it. This proves (B.25.1).

For the quartic put $L_n=\mathbb Q(E[\ell^n])$ and
let $\Gamma_n$ be its actual representation image.
The closed CM image contains $-I$.
Indeed the residual Cartan in B.23 contains $-I$.
A lift in its commutative CM unit group has eigenvalues
in $-1+\ell\mathcal O$; its successive
$\ell^k$-th powers converge to $-I$.
The image is closed, so the limit belongs to it.
For a cocycle on $\Gamma_n$ with coefficients $V/\ell$,
commuting an element with this central scalar proves
that the cocycle is a coboundary, since $-2$ is a unit.
This is the direct cocycle proof of the scalar
vanishing $H^1(\Gamma_n,V/\ell)=0$.

The reduction $h_1$ is nonzero by B.23.
Inflation--restriction implies that its restriction
to $\ker(G_S\to\Gamma_n)$ is nonzero.
The residual representation is irreducible:
at three the nonsplit Cartan already acts irreducibly;
at five the two distinct Cartan eigenlines are
exchanged by complex conjugation.
Thus the image of $h_n$ on that kernel is onto modulo
$\ell$, and Nakayama's lemma gives all of $V/\ell^n$.
Consequently the actual finite quotient has image
\[
 g\longmapsto(h_n(g),\rho_n(g)):
 \qquad (V/\ell^n)\rtimes\Gamma_n .
                                                               \tag{B.25.4}
\]

The $\ell$-Sylow subgroup of $\Gamma_n$ lies in the
CM principal units modulo $\ell^n$ and has exponent
at most $\ell^{n-1}$. Hence $\mu_{\ell^{n+1}}$ is
not contained in $L_n$: its cyclotomic group would
have a quotient of order $\ell^n$.
Adjoining $F=\mathbb Q(\zeta_8)$ preserves this fact,
since its degree is prime to $\ell$.
Over $L_nF$, the Kummer extension generated by
$2^{1/\ell^n}$ has cyclic $\ell$-kernel with
conjugation character $\chi_\ell$, whereas the
next cyclotomic degree-$\ell$ extension has trivial
conjugation character. These characters are
different modulo $\ell$. Therefore one can fix
$L_nF$ and $2^{1/\ell^n}$ while acting nontrivially
on $\mu_{\ell^{n+1}}$.

The additional extension just described has
$\ell$-primary composition factors with characters
$1$ or $\chi_\ell$, and a prime-to-$\ell$ part.
Its common quotient with the translation extension
in (B.25.4) is trivial: every nonzero quotient
of the latter has the irreducible two-dimensional
$V/\ell$ as a composition factor.
Thus the preceding automorphism and a primitive
translation in (B.25.4) can be prescribed
simultaneously. All fields and their conjugation
actions are actual normal extensions over $\mathbb Q$.
Chebotarev applied to that conjugacy class
\citep[VIII, Theorem 7.4]{BMilneCFT2020}
gives infinitely many primes satisfying the first
five conditions in (B.25.2).
Since $F_E-I$ is divisible by $\ell^n$ while
$h(F_p)$ is primitive, its inverse has denominator
at least $\ell^n$. Hence $m_p\geq n$.

Pure compact inputs lift by B.24.
For (B.25.3), every shifted peripheral cup is a
scalar multiple of $\kappa_\ell(2)\cup\kappa_\ell(p)$:
the two self-cups vanish since two is a unit.
At the primes over $p$, the residue field is
$\mathbb F_p$, and its entire $\ell$-primary
tame symbol is zero by the power and valuation
conditions in (B.25.2).
At the prime over two the residue multiplicative
group is trivial; there are no other supports.
Thus all peripheral tests hold, so this toric
input also lifts. Its mixed inertia vector at
$p$ is primitive. The order bound $m_p\geq n$
then violates (B.24.17) when $n>t$.
This proves the asserted failure of additivity.
\end{proof}

\begin{statement}{B.26}{The least positive derived length for the quintic is one}{Proved}
Every Selmer section of the full profinite geometric
abelianization of $U_5$ belongs to the image of a rational
cusp decomposition group. Consequently the least positive
integer in the quintic derived-length question is $k=1$.
The maximal-prosolvable Selmer set is likewise cuspidal,
by the separate finite-support argument B.1 and A.50--A.52.
\end{statement}
\begin{proof}
We first record the exact global-local fact needed for
the torus. For every number field $L$ the map
\[
 H^1(L,\widehat{\mathbb Z}(1))
      \longrightarrow\prod_{w\ {\rm finite}}
                H^1(L_w,\widehat{\mathbb Z}(1))
                                                               \tag{B.26.1}
\]
is injective. Work prime by prime in its coefficient.
For an odd prime $\ell$, the local-power theorem
\citep[VIII, Theorem 1.4]{BMilneCFT2020} says that a
class in $L^\times/L^{\times\ell^n}$ which is zero
at all but finitely many places is zero globally.
The condition concerning the cyclic two-primary
cyclotomic extension in that theorem is empty here.
Finite roots-of-unity groups are Mittag--Leffler,
so $H^1(L,\mathbb Z_\ell(1))$ is the inverse limit
of these finite-coefficient Kummer groups.
This proves the assertion at odd $\ell$.

At $\ell=2$, first restrict a locally zero class
$z$ to $L(i)$. Every extension
$L(i,\mu_{2^n})/L(i)$ is cyclic, being a quotient
of a subgroup of the cyclic group
$\operatorname{Gal}(\mathbb Q(\mu_{2^n})/\mathbb Q(i))$.
The same theorem therefore makes every reduction
of this restriction zero. If $i\notin L$,
restriction followed by corestriction gives $2z=0$;
if $i\in L$, it already gives $z=0$.
The exact sequence
\[
 0\longrightarrow\mathbb Z_2(1)
 \longrightarrow\mathbb Q_2(1)
 \longrightarrow(\mathbb Q_2/\mathbb Z_2)(1)
 \longrightarrow0
\]
identifies the torsion subgroup of
$H^1(L,\mathbb Z_2(1))$ with
$\mu_{2^\infty}(L)$: $H^0(L,\mathbb Q_2(1))=0$,
and $H^1(L,\mathbb Q_2(1))$ is a vector space.
This finite torsion embeds into the corresponding
local group at any dyadic place, by the identical
sequence there and the inclusion of roots of unity.
Thus the locally zero $z$ is zero. This proves
(B.26.1), without the false assertion that every
finite two-primary Kummer map over an arbitrary
number field is injective.

Now let $X=X_5$ and $J=J_X$. The three cyclic
maps to $C:y^2+y=x^5$ in A.53 induce an isogeny
$J\sim_{\mathbb Q}J_C^3$. A.50--A.52 establish
rank zero and finite Tate--Shafarevich group for
$J_C$, with the precise published arithmetic
input stated there. These properties are invariant
under an isogeny, so they hold for $J$.
Fix a rational cusp $O$ and use its unit tangent
as section origin. For
$V=TJ_{X,D}$ there is the full profinite sequence
\[
 0\longrightarrow W\longrightarrow V\longrightarrow TJ
     \longrightarrow0,\qquad
 W=\widehat{\mathbb Z}(1)[D]/
       \widehat{\mathbb Z}(1)\sum_{d\in D}d .
                                                               \tag{B.26.2}
\]
The compact projection of a Selmer section lies in
the full Tate-module Selmer group of $J$.
Rank zero and finite $\Sha$ identify this group
with the finite group $J(\mathbb Q)$.
Let $t$ be the corresponding point.
At every finite $p$ its local value is the Kummer
class of $[P_p-O]$, where $P_p\in X(\mathbb Q_p)$
is the underlying local point or cusp.
The full local Kummer map is injective because
$J(\mathbb Q_p)$ is a compact $p$-adic Lie group
with trivial intersection of all multiplication
subgroups. Therefore $[P_p-O]=t$.
The Abel--Jacobi map is a closed embedding.
Since its fibre over the rational point $t$ is
nonempty at a completion, it is its single
rational point $c\in X(\mathbb Q)$.
Thus every $P_p=c$. A.52 shows that every such
rational point of $X$ is a cusp.

Subtract the unit tangent section at $c$.
The resulting class has zero compact projection,
so it comes from $z\in H^1(\mathbb Q,W)$.
The ambiguity from $H^0(\mathbb Q,TJ)$ is zero.
Locally the class belongs to the inertia line
$I_c=\widehat{\mathbb Z}(1)e_c$.
Moreover $H^0(\mathbb Q_p,TJ)=0$, since
$J(\mathbb Q_p)$ has finite torsion. Hence
$H^1(\mathbb Q_p,W)\to H^1(\mathbb Q_p,V)$
is injective, and the image of $z$ in
$H^1(\mathbb Q,W/I_c)$ is locally zero.

Choose another rational cusp $d\ne c$.
Eliminating its coordinate in the diagonal
relation gives the integral permutation isomorphism
\[
 W/I_c\simeq
 \widehat{\mathbb Z}(1)[D\setminus\{c,d\}]
 \simeq\widehat{\mathbb Z}(1)\oplus
 \left(\operatorname{Ind}_{G_{\mathbb Q(\zeta_5)}}^{G_{\mathbb Q}}
                   \widehat{\mathbb Z}(1)\right)^3 .
                                                               \tag{B.26.3}
\]
Shapiro's lemma, locally as well as globally,
and (B.26.1) show that this locally zero class
is zero. Exactness then places $z$ in
the image of $H^1(\mathbb Q,I_c)$.
Every such cocycle is a section of the actual
split cusp decomposition group at $c$.
Its image in the abelian arithmetic quotient is
the original section, proving the assertion.

As in B.5, this least-positive-length conclusion
is not a monotonicity theorem in the derived
length. It does not classify arbitrary
geometrically pro-$5$ section families, and it
does not infer higher quotient cuspidality
solely from the abelian statement.
\end{proof}

\begin{statement}{B.27}{Every unramified three- and five-primary family has exactly three Selmer packets}{Proved}
Let $\ell\in\{3,5\}$, $n=\ell^r$ with $r\geq1$, and
let $U_n$ be the plus Fermat curve with its coordinate
cusps removed. Fix the rational unit tangent at
$[0:-1:1]$ and its genuine arithmetic action on
\[
 N_n=\pi_1(U_{n,\overline{\mathbb Q}})^{(\ell)}.
\]
Let $J$ be any closed normal, $G_{\mathbb Q}$-stable
subgroup with $J\subseteq N_n'$, and put $Q=N_n/J$.
For an actual section of $Q\rtimes G_{\mathbb Q}$,
assume that its geometric cocycle factors through
$G_{\mathbb Q,S}$, where $S=\{\ell,\infty\}$.
This is the precise unramified-outside-$\ell$ hypothesis.
Then the following conditions are equivalent:

\begin{enumerate}
\item Its localization at $\ell$ is the image of a
local point section or a local cusp section.
\item It is Selmer in this quotient at every place.
\item It is the image of a global section of the
decomposition group of one of the three rational cusps.
\end{enumerate}
The section classes satisfying these conditions are
the disjoint union of three copies of $\mathbb Z_\ell$.
Relative to the unit tangent at the chosen cusp $c$,
the parameter $a\in\mathbb Z_\ell$ is precisely
$a\kappa_\ell(\ell)$ in its inertia
$I_c\simeq\mathbb Z_\ell(1)$.
The assertion includes $J=1$, every positive-class
intrinsic pro-$\ell$ nilpotent quotient, and every
positive-length intrinsic pro-$\ell$ derived quotient.

Consequently, for every actual lift constructed in
LM6 E.16 and LM7 E.34, Selmerness is equivalent to
being one of these global cusp images. Every other
member fails local geometricity already at $\ell$,
including the noncuspidal family of LM7\_PUSH B.2.
Together with B.6 this determines the Selmer status
of all families listed in R9. No statement about
arbitrary sections ramified away from $\ell$ is
included here.
\end{statement}
\begin{proof}
\emph{The coefficient-prime decomposition group is onto.}
Put $K=\mathbb Q(\zeta_\ell)$ and let $\lambda$
be its unique prime above $\ell$.
For $\ell=3$ the Eisenstein norm is Euclidean.
For $\ell=5$ Minkowski's ideal-class bound is
\[
 \frac{4!}{4^4}\left(\frac4\pi\right)^2\sqrt{125}<2,
\]
so every ideal class has an ideal of norm one.
Thus both fields have class number one.
The identity
$\ell=\text{unit}\cdot(1-\zeta_\ell)^{\ell-1}$
gives the unique prime $\lambda$.
The fresh program \texttt{b\_unramified\_selmer.gp}
also certifies both class numbers and ramification
decompositions exactly.

Let $P=G_{K,\{\lambda\}}^{(\ell)}$.
The image $D$ of $G_{K_\lambda}$ in $P$ is closed.
If $D\ne P$, some finite $\ell$-group quotient
has a proper image of $D$. Place that image in
a maximal subgroup, necessarily normal of index
$\ell$. The resulting cyclic extension of $K$
is unramified outside $\lambda$, and $\lambda$
splits completely in it. It is therefore everywhere
unramified. Class field theory and class number
one make such an extension impossible.
Hence $G_{K_\lambda}\to P$ is surjective.

Write $H=G_{K,\{\lambda\}}\subset G_{\mathbb Q,S}$
and define the actual relative arithmetic quotient
\[
 \Gamma=G_{\mathbb Q,S}/
         \ker(H\longrightarrow P),\qquad
 1\longrightarrow P\longrightarrow\Gamma
       \longrightarrow\operatorname{Gal}(K/\mathbb Q)
       \longrightarrow1.
                                                               \tag{B.27.1}
\]
The kernel being divided out is characteristic in
$H$ and hence normal in $G_{\mathbb Q,S}$.
The decomposition group $G_{\mathbb Q_\ell}$
surjects onto the last finite group, since
$K/\mathbb Q$ is totally ramified at $\ell$.
Its intersection with $H$ maps onto $P$, as just
proved. Therefore
\[
             G_{\mathbb Q_\ell}\longrightarrow\Gamma
                  \quad\text{is surjective}.                  \tag{B.27.2}
\]
No freeness assertion about $P$ is needed for this step.

\emph{Factorization of the genuine action and every cocycle graph.}
Let $F_2$ denote the free geometric pro-$\ell$
tripod group with its standard unit tangent.
The actual power cover identifies
\[
 N_n=\ker\left(F_2\longrightarrow(\mathbb Z/n)^2\right).
                                                               \tag{B.27.3}
\]
At the chosen cusp one may write
$X/Z=-s$, $t=s^n$ and $Y/Z=-(1-s^n)^{1/n}$.
The leading constants have coherent rational
$\ell$-power roots, so coefficientwise action on
these Puiseux fibres gives the genuine restriction
of the tripod action in (B.27.3).
The pair and its unit tangent have good reduction
outside $\ell$; prime-to-residue-characteristic
specialization gives the unramified based action,
as in the precise argument of B.24.

On $F_2/\Phi(F_2)\simeq\mathbb F_\ell^2$ the
action is the cyclotomic scalar: these coordinates
are the two peripheral abelian characters.
It is the identity over $K$.
The kernel of
$\operatorname{Aut}(F_2)\to
 \operatorname{GL}(F_2/\Phi(F_2))$ is pro-$\ell$.
Indeed the characteristic lower exponent-$\ell$
quotients are finite $\ell$-groups, and the
automorphisms inducing the identity on their
Frattini quotients form $\ell$-groups by the
finite basis theorem. The successive kernels
embed in homomorphism groups of their finite
$\mathbb F_\ell$-graded pieces.
Taking their inverse limit proves the assertion.
Restrict the compact arithmetic subgroup
$\rho(H)$ to the characteristic open subgroup
(B.27.3). It preserves $J$ by the assumed
Galois stability, so its action descends to $Q$.
Thus the actual action image of $H$ on $Q$
is pro-$\ell$. No assertion that every IA
automorphism preserves $J$ is required.

Let $z:G_{\mathbb Q,S}\to Q$ be any continuous
section cocycle under consideration.
Its graph homomorphism on $H$ has target
\[
 Q\rtimes\rho_Q(H),
 \qquad h\longmapsto(z(h),\rho_Q(h)),
\]
a pro-$\ell$ group. It therefore kills the
defining kernel of $H\to P$. Both $z$ and
the action factor through the same actual
quotient $\Gamma$ in (B.27.1).
It is this graph in $Q\rtimes\operatorname{Aut}(Q)$
which factors through $\Gamma$; the section into
$Q\rtimes G_{\mathbb Q}$ retains its identity
arithmetic projection and does not so factor.
This proves the graph assertion for every cocycle
in the stated unramified space, including
the particular lifts used in earlier returns.
It also shows that restriction of these section
classes to $G_{\mathbb Q_\ell}$ is injective:
after a local geometric conjugation, equality
on that group is equality on $\Gamma$ by
(B.27.2), hence equality globally.

\emph{A local geometric section at $\ell$ has a rational
cusp centre.}
First consider $X_\ell$ and write
$A_\ell=\operatorname{Jac}(X_\ell)$.
A.10 and A.14 give rank zero and finite $\Sha$
for $A_3$; A.50--A.53 give the same conclusions
for $A_5$, with the precise Known arithmetic
input stated there.
The local groups $A_\ell(\mathbb Q_\ell)$
are pro-$\ell$.
For five this is A.54.
For three the minimal model
$y^2+y=x^3-7$ has discriminant $-3^9$,
Kodaira type $\mathrm{IV}^*$ and Tamagawa number
three. The exact local reduction is
$[3,-4,[1,0,0,0],3]$ in the convention of
\citet[\texttt{elllocalred}]{BPARIElliptic2026};
it is recomputed in \texttt{b\_unramified\_selmer.gp}.
Its nonsingular reduction has three points,
and its formal group is pro-three.
The reduction and component exact sequences
therefore make $A_3(\mathbb Q_3)$ pro-three.
In particular the single-coefficient Kummer map
\[
 A_\ell(\mathbb Q_\ell)\longrightarrow
           H^1(\mathbb Q_\ell,T_\ell A_\ell)
                                                               \tag{B.27.4}
\]
is injective. Notice that this is stronger than
the full profinite Kummer injection used in B.26.

The cover $U_n\to U_\ell$ raises each projective
coordinate to $n/\ell$; it extends to the proper
curves and sends the unit tangent to the unit
tangent, with inertia multiplied by $n/\ell$.
Since $J\subseteq N_n'$, the quotient retains
$N_n^{\rm ab}$ and thus the composite compact
map to $T_\ell A_\ell$.
The projected cohomology class is unramified at
every finite $p\ne\ell$.
At such a prime $A_\ell$ has good reduction,
and unramified $T_\ell A_\ell$ classes are
exactly the Kummer image of
$A_\ell(\mathbb F_p)[\ell^\infty]$.
At the real place the odd-primary $H^1$ is zero.
The assumed local point-or-cusp image at $\ell$
places the remaining localization in the Kummer
image there. Thus the compact class belongs
to the global $\ell$-primary Tate Selmer group.

Rank zero and finite $\Sha$ identify that group
with $A_\ell(\mathbb Q)^{\wedge_\ell}$.
The finite group $A_\ell(\mathbb Q)$ injects
into the pro-$\ell$ local group, so it is
already $\ell$-primary; its completion is itself.
Let $a$ be the resulting rational point.
If $P_\ell$ is the compact underlying local
point on $X_\ell$, (B.27.4) gives
$[P_\ell-O]=a$ in $A_\ell(\mathbb Q_\ell)$.
The Abel--Jacobi embedding then makes
$P_\ell$ the corresponding rational point
of $X_\ell$. A.10 and A.52 say that every
such rational point is a coordinate cusp.

An interior point of $U_n$ cannot map to a cusp
of $X_\ell$, because none of its three coordinates
vanishes. Therefore the local section on $U_n$
must be the image of a local cusp section.
Moreover every $\mathbb Q_\ell$-rational cusp
of $X_n$ is already one of the three rational
cusps over $\mathbb Q$: a coordinate ratio is
an $n$-th root of $-1$, and
$\mathbb Q_\ell$ contains no nontrivial
$\ell$-power root of unity.
This proves that the actual local branch is
a single rational cusp $c$ of $U_n$.

\emph{From this local branch to its global packet.}
Rebase the split arithmetic extension at the
rational unit tangent of $c$.
That tangent is also unramified outside $\ell$,
and the difference from the initial tangent
is a cocycle on $G_{\mathbb Q,S}$.
Normality and Galois stability of $J$ make
this change of action well defined on $Q$.
Use the geometric conjugator which aligns
the localized section with the local cusp
section as a single global conjugator in $Q$.
This preserves unramifiedness, because the
based action on inertia outside $\ell$ is
trivial. The resulting cocycle $z_c$ has
\[
             z_c(G_{\mathbb Q_\ell})\subset I_c
                        \subset Q .
                                                               \tag{B.27.5}
\]
Here $I_c$ is the actual inertia image, stable
under the cusp-based arithmetic action.
Because $J\subseteq N_n'$, its map
$\mathbb Z_\ell(1)\to Q$ is injective:
the peripheral generator is primitive in
$N_n^{\rm ab}$.
The original cocycle and the cusp-tangent change
both have graphs factoring through $\Gamma$
by the preceding argument. Their product and
the corresponding inner change of action do too.
Thus the rebased graph still factors through
$\Gamma$. Surjectivity (B.27.2) now implies
$z_c(G_{\mathbb Q,S})\subset I_c$ directly:
every value is the value of a local element.
No normalizer theorem or nonabelian
restriction/corestriction argument is invoked.

Kummer theory gives
\[
 H^1(G_{\mathbb Q,S},I_c)
   =\mathbb Z_\ell\,\kappa_\ell(\ell).
                                                               \tag{B.27.6}
\]
Indeed unramifiedness kills all completed
prime-valuation coordinates except that at
$\ell$, and $\{\pm1\}$ has zero odd-primary
Kummer completion.
Thus (1) implies precisely the global cusp
description in (3), with the stated parameter.
Each such parameter comes from a global full
cusp section: choose this $\ell$-component
and zero at every other coefficient prime
in $H^1(\mathbb Q,\widehat{\mathbb Z}(1))$.
It is therefore Selmer in the stipulated
quotient sense. This proves (3)$\Rightarrow$(2),
and (2)$\Rightarrow$(1) is immediate.

Finally the three packets and their parameters
are distinct as geometric-conjugacy classes.
Within one packet, $H^0(\mathbb Q,V_n/I_c)=0$,
where $V_n=N_n^{\rm ab}$: at a good Frobenius
its eigenvalues have absolute values $\sqrt p$
or $p$, hence none is one.
Consequently $H^1(\mathbb Q,I_c)\to
H^1(\mathbb Q,V_n)$ is injective.
Different centres are distinguished by their
compact image on $X_\ell$: the three rational
cusp points have distinct Abel--Jacobi classes
in the $\ell$-primary group $A_\ell(\mathbb Q)$.
This proves the disjoint three-packet
parametrization.

LM6 E.16 and LM7 E.34 use exactly this
unramified arithmetic quotient in their
construction, as their defining graph maps
show. The same is true of the lift used in
LM7\_PUSH B.2; its toric vector
$(\kappa_3(3),2\kappa_3(3))$ lies on none
of the three cusp lines.
Thus every noncuspidal chosen lift through this
unramified source fails condition (1), even when
its abelian class itself is cuspidal. The
hypothesis concerns the whole cocycle, and
does not classify a different lift with added
ramification merely because its abelian class
is unramified.
The separately considered full profinite
metabelian families have exactly the Selmer
classification proved in B.6.
\end{proof}

\begin{statement}{B.28}{Every nonempty intrinsic class-three fibre is evaluated}{Proved}
Let $N$ be the actual geometric pro-three group of $U_3$,
with the unit tangent at $C_0=[1:-1:0]$ for $q_0=Z/X$
as in B.24, and put
$N_i=N/\gamma_{i+1}N$ and $A=\gamma_3N/\gamma_4N$.
For every input class, ramified inputs included, every
nonempty fibre of
\[
 H^1(\mathbb Q,N_3)\longrightarrow H^1(\mathbb Q,N_2)
                                                               \tag{B.28.1}
\]
is a torsor under the explicitly evaluated group
\[
 \mathbb Z_3^{166}\oplus(\mathbb Z/3)^7
 \oplus(\mathbb Z/9)^{56}\oplus(\mathbb Z/27)^{25}
 \oplus(\mathbb Z/81)^3 .                                      \tag{B.28.2}
\]
Every whole cocycle unramified outside three has a nonempty
fibre. This determines all nonempty fibres in prompt B.5
without using a value of the Ross class. It does not by
itself decide the image for arbitrary ramified inputs.
\end{statement}
\begin{proof}
The sequence $1\to A\to N_3\to N_2\to1$ is central,
and $A=\operatorname{Lie}_3(V)$ for $V=N^{\rm ab}$.
Its arithmetic action is the third Lie functor of the
actual degree-one action; no degree-two or degree-three
arithmetic correction changes that induced action.
The rank of $A$ is $(10^3-10)/3=330$.

Fix a lift of an input section. Every other lift differs
from it by an $A$-valued cocycle. There is no further
quotient by a stabilizer in $N_2$: if an element were fixed
by the section-twisted arithmetic action, its image in $V$
would be fixed by the original arithmetic action.
Good Frobenius weights give $V^{G_{\mathbb Q}}=0$.
It would therefore lie in $\Lambda^2V$, on which the inner
twist is trivial, and the same weights give
$(\Lambda^2V)^{G_{\mathbb Q}}=0$.
Thus that element is one. Two lifted cocycles are
geometrically conjugate exactly when their difference is
an $A$-coboundary. The fibre is consequently a torsor
under $H^1(\mathbb Q,A)$.

At $p\ne3$, the action on $A$ is unramified. An
$A$-cocycle sends wild inertia to zero and sends a
three-primary tame generator to a vector $v$ satisfying
$\Phi_{A,p}v=pv$. The rational eigenvalues of
$\Phi_{A,p}$ have complex absolute values
$p^{3/2},p^2,p^{5/2},p^3$, so $v=0$.
Every such cocycle is therefore unramified outside three,
and
\[
 H^1(\mathbb Q,A)=H^1(G_{\mathbb Q,\{3,\infty\}},A).
                                                               \tag{B.28.3}
\]
The finite-base $H^2$ vanishing in B.7, with its finite
$H^1$ inverse systems, also shows that every unramified
input has a lift.

Here is a complete integral evaluation of (B.28.3).
Use $U,R_3$ from (B.24.2), and set
\[
 \begin{gathered}
 C_E=\begin{pmatrix}1&-1\\0&-1\end{pmatrix},\qquad
 C_T=-P_{(A_1\ A_2)(B_1\ B_2)(C_1\ C_2)},\\
 T_V=\begin{pmatrix}
 I_2+3R_3&0\\
 U(I_2+3R_3-7I_2)/3&7I_8
 \end{pmatrix},\qquad
 C_V=\begin{pmatrix}
 C_E&0\\(UC_E-C_TU)/3&C_T
 \end{pmatrix}.
 \end{gathered}                                                 \tag{B.28.4}
\]
The divided blocks are integral by the actual cusp
three-torsion coordinates and inverse Weil-pairing
comparison proved in B.24. The matrix $C_V$ is complex
conjugation. The matrix $T_V$ is the Frobenius at the
chosen prime of $K=\mathbb Q(\zeta_3)$ over seven:
its primary CM scalar is $1+3\zeta_3$.
The program below counts $E_{27}(\mathbb F_7)$ directly
and obtains nine; the norm, trace and primary congruence
therefore identify this scalar up to conjugation,
which is fixed by choosing that prime.

The maximal pro-three group $P$ of $K$ unramified outside
three is free of rank two. Indeed class number one and
Kummer theory give $\dim H^1(P,\mathbb F_3)=2$;
the finite-base $H^2$ vanishing gives
$H^2(P,\mathbb F_3)=0$. Indeed, in the inflation sequence an
invariant character of the kernel of the maximal
pro-three quotient would give a larger pro-three quotient,
so the adjacent kernel-character group is zero.
Filter finite abelian three-primary kernels by trivial
simple modules and then induct on the centre of a finite
three-group kernel. This solves every finite pro-three
embedding problem for $P$. A minimal free presentation
therefore splits; its rank-two free source and the
Frattini criterion make the induced idempotent the
identity. This proves freeness.
The actual CM image lies in $1+3\mathbb Z_3[\zeta_3]$.
The logarithms of the two primary scalars
$1+3\zeta_3$ and $1+3\zeta_3^2$, divided by three,
reduce to $\zeta_3,\zeta_3^2$ modulo three, a basis
of the rank-two additive lattice. Thus the two actual
Frobenii $\tau,c\tau c$ form a Frattini basis of $P$,
and hence a free pro-three basis.
The relative arithmetic group is $P\rtimes\langle c\rangle$.
The graph argument of B.27 shows that every cocycle in
(B.28.3) factors this relative group.

Let $T=\operatorname{Lie}_3(T_V)$ and
$C=\operatorname{Lie}_3(C_V)$.
Since two is invertible, every cocycle has a representative
vanishing at $c$. Such a representative is freely determined
by its value at $\tau$; its value at $c\tau c$ is $C$
times that value. The coboundaries preserving this
normalization are $(T-I)A^c$. Therefore
\[
 H^1(\mathbb Q,A)
   =\operatorname{coker}\bigl((T-I)(I+C):A\to A\bigr).
                                                               \tag{B.28.5}
\]
The trace of $C_V$ is $-2$. The degree-three Witt
trace formula gives
$\operatorname{tr}C=((-2)^3-(-2))/3=-2$.
Hence $A^c$ has rank $164$.
The good-Frobenius weights make $T-I$ invertible
rationally, so the matrix in (B.28.5) has exact rank
$164$ and free cokernel rank $166$.

The fresh program \texttt{b\_class3\_fibres.py} constructs
all $330$ Hall words $[[v_i,v_j],v_k]$ with $i>j$ and
$k\ge j$, and both Lie-action matrices from (B.28.4).
It uses the integral Jacobi identity for every reduction.
Unimodular three-adic elimination at modulus $3^{12}$
gives the $164$ nonzero Smith valuations
\[
              0^{73},\quad1^7,\quad2^{56},\quad3^{25},\quad4^3.
                                                               \tag{B.28.6}
\]
Every pivot is strictly below the working precision.
Together with the independently proved exact rank, this
certifies all integral elementary divisors, rather than
only a truncated matrix rank. It proves (B.28.2).
The script reads neither inherited matrices nor old output.
\end{proof}

\subsection{A full pro-three counterexample ramified at seven}

The implication in task B.10(b) is false, already for the actual
geometric pro-three group.  The following finite-set arithmetic
calculation is what permits us to prescribe the local cusp packet
without replacing an actual Frobenius word by an independent symbol.

\begin{statement}{B.29}{The two local factors above seven}{Proved}
Put
\[
 K=\mathbb Q(\zeta),\quad \zeta^2+\zeta+1=0,\quad
 \lambda=1-\zeta,\quad \pi=2-\zeta,\quad
 \bar\pi=3+\zeta,
\]
and let $S_K$ consist of the three primes
$(\lambda),(\pi),(\bar\pi)$.  Thus $\pi\bar\pi=7$.
Write
\[
 P_S=G_{K,S_K}^{(3)},\qquad
 L_1=G_{K_{(\pi)}}^{(3)},\qquad
 L_2=G_{K_{(\bar\pi)}}^{(3)}.
\]
After choosing prolongations, the actual continuous local maps
induce an isomorphism of pro-three groups
\[
             L_1*_3 L_2\ \xrightarrow{\ \sim\ }\ P_S.       \tag{B.29.1}
\]
Here each $L_i$ is the maximal pro-three group of $\mathbb Q_7$,
with its tame presentation
\[
 L_i=\langle t_i,F_i\mid F_it_iF_i^{-1}=t_i^7\rangle_{\rm pro\text{-}3}.
                                                               \tag{B.29.2}
\]
In particular the actual local maps are injective; this is a
conclusion, not an assumption of the construction.
Choose the two prolongations to be exchanged by a fixed complex
conjugation $c\in G_{\mathbb Q}$.  The actual quotient
\[
 \Gamma_S=
 G_{\mathbb Q,\{3,7,\infty\}}/
 \ker(G_{K,S_K}\longrightarrow P_S)
       =P_S\rtimes\langle c\rangle                             \tag{B.29.3}
\]
has $c$ exchanging the two factors of (B.29.1).
\end{statement}
\begin{proof}
Let $X=\operatorname{Spec}\mathcal O_K[1/21]$.
The Eisenstein integer ring is norm Euclidean, hence a principal
ideal domain.  Its unit group is $\mu_6$, and its localized unit
group is generated by $\mu_6,\lambda,\pi,\bar\pi$.
Kummer theory consequently identifies
\[
 H^1(P_S,\mathbb F_3)
   =H^1_{\rm et}(X,\mu_3)
   =\langle\zeta,\lambda,\pi,\bar\pi\rangle_{\mathbb F_3}.
                                                               \tag{B.29.4}
\]
We identify $\mathbb F_3$ with $\mu_3$ by $1\mapsto\zeta$.
The first equality only uses degree-one comparison: a cubic
Kummer torsor is itself a three-cover.  Each local group $L_i$
has two-dimensional $H^1$, since
$\mathbb Q_7^\times/(\mathbb Q_7^\times)^3$ has one valuation
coordinate and one residue-unit coordinate.

At $(\pi)$ the reduction of $\zeta$ is $2$, and at $(\bar\pi)$
it is $4$.  The reductions of $\lambda$ are respectively $6$
and $4$.  On arithmetic Frobenius, the Kummer character of a
unit $u$ is determined by $\bar u^{(7-1)/3}=\bar u^2$.
Thus the two Frobenius rows on $\zeta,\lambda$ are
\[
                   \begin{pmatrix}2&0\\2&2\end{pmatrix}.
                                                               \tag{B.29.5}
\]
The inertia rows on $\pi,\bar\pi$, after independently choosing
positive tame generators, form the identity matrix.  They vanish
on $\zeta,\lambda$.  Therefore the full restriction matrix has
the block form
\[
 \begin{pmatrix}
 2&0&*&*\\2&2&*&*\\0&0&1&0\\0&0&0&1
 \end{pmatrix},\qquad \det=1\quad\text{in }\mathbb F_3.          \tag{B.29.6}
\]
The unspecified entries retain the actual choices of Frobenius;
they have no effect on this determinant.  By the universal
property of the free pro-three product, (B.29.6) says that
restriction from $P_S$ to $B=L_1*_3L_2$ is an isomorphism in
$H^1(-,\mathbb F_3)$.

We also have an injection
\[
 H^2(P_S,\mathbb F_3)\lhook\joinrel\longrightarrow
 H^2_{\rm et}(X,\mu_3).                                      \tag{B.29.7}
\]
This is the general degree-two injection preceding
\citet[Proposition 2.1, author version p.~4]{CBSchmidt2007}:
the first cohomology of the universal pro-three cover vanishes,
and the Hochschild--Serre five-term sequence gives the injection.
For clarity, a cyclic three-torsor on a finite Galois three-cover
is killed on a further Galois three-cover: take its finitely
many conjugates and their compositum.  Its Galois group over
the original base is an extension of a subgroup of a product
of cyclic groups of order three by a three-group.  This proves
the vanishing used in that five-term argument.  We do not need
a blanket comparison in all cohomological degrees.

Since $\operatorname{Pic}(X)=0$, the Kummer sequence gives
$H^2_{\rm et}(X,\mu_3)=\operatorname{Br}(X)[3]$.
The Brauer group of the regular integral curve $X$ injects
into $\operatorname{Br}(K)$
\citep[Theorem 3.5.4]{CBColliotTheleneSkorobogatov2019}.
An element of this Brauer group has zero local invariant outside
$S_K$: at a finite point of $X$ it extends over the completed
local integer ring, whose Brauer group is zero because that
ring is henselian with finite residue field
\citep[Theorem 3.4.2]{CBColliotTheleneSkorobogatov2019}.
There are no real places.  Global Brauer reciprocity
\citep[VIII, Theorem 4.2]{BMilneCFT2020} therefore implies that
restriction to the two primes above seven is injective on
$\operatorname{Br}(X)[3]$.  Indeed, if those two invariants
vanish, the sole remaining invariant at $(\lambda)$ is zero
by their sum, and all local invariants then vanish.
By naturality of (B.29.7), a class in $H^2(P_S,\mathbb F_3)$
whose restrictions to $L_1$ and $L_2$ are zero is zero.
In particular
\[
                H^2(P_S,\mathbb F_3)\longrightarrow
                       H^2(B,\mathbb F_3)
                      \quad\text{is injective}.              \tag{B.29.8}
\]
This argument does not assume in advance that either local
group embeds in $P_S$.

Let $f:B\to P_S$ be the actual local map.  Its image is compact
and hence closed.  If this image were proper, its image would be
proper in a finite three-group quotient of $P_S$ and would lie
in a maximal subgroup of index three.  The corresponding
nonzero $H^1(P_S,\mathbb F_3)$ character would restrict to zero
on $B$, contradicting (B.29.6).  Thus $f$ is surjective.
For $R=\ker f$, its five-term sequence and (B.29.6)--(B.29.8)
give
\[
                  H^1(R,\mathbb F_3)^{B}=0.                  \tag{B.29.9}
\]
If $R$ were nontrivial, choose a finite three-group quotient
$B_0$ of $B$ in which its image $R_0$ is nontrivial.
The nonzero module $R_0/\Phi(R_0)$ has nonzero coinvariants
under $B_0$: the augmentation ideal of $\mathbb F_3[B_0]$ is
nilpotent.  A nonzero linear functional on those coinvariants
gives a nonzero $B$-invariant character of $R$, contrary to
(B.29.9).  This proves (B.29.1).

For (B.29.2), wild inertia is pro-seven and hence dies in every
pro-three quotient.  Tame inertia has three-primary part
$\mathbb Z_3(1)$ and the unramified quotient has three-primary
part $\mathbb Z_3$.  Since $7\equiv1\pmod3$, its action by
$7$ is continuous on this quotient; the tame presentation
therefore gives exactly (B.29.2).
Finally $K/\mathbb Q$ is Galois and is ramified only at three,
so $G_{K,S_K}$ is the actual index-two subgroup of the rational
restricted-ramification group in (B.29.3).
Its maximal-pro-three kernel is characteristic and hence normal
in that group.  The image of the chosen complex conjugation
has order two and splits the quotient.  Taking conjugate
prolongations makes its exchange of local factors literal.
There is no independent prescription of a higher Frobenius
word anywhere in the argument.
\end{proof}

\begin{statement}{B.30}{Ramification at seven does not force a rational cusp}{Refuted}
Let
\[
 U_3=\{X^3+Y^3+Z^3=0\}\setminus\{XYZ=0\},\qquad
 N=\pi_1(U_{3,\overline{\mathbb Q}})^{(3)}.
\]
Use the genuine arithmetic action based at the rational unit
tangent at $O=[0:-1:1]$.  The assertion that every section
ramified at a prime different from three belongs to a global
rational-cusp packet is false.  A counterexample is supplied
by each nonzero member of the following family.
There is a family, indexed by
$a\in\mathbb Z_3$, of actual continuous sections
\[
                    s_a:G_{\mathbb Q}\longrightarrow
                                  N\rtimes G_{\mathbb Q}
\]
with pairwise distinct geometric conjugacy classes, unramified
outside $\{3,7\}$.  At seven, $s_a$ is the actual local cusp
section centred at
\[
                         P=[0:-\zeta:1]
\]
in the embedding $\zeta\equiv2\pmod7$, with tangent twisted by
$a\kappa_3(7)$.  For every $a\ne0$ it is ramified at seven
and is not the image of a global rational cusp section.

The same conclusions hold after projecting to every quotient
$N/J$ for which $J$ is closed, normal, $G_{\mathbb Q}$-stable,
and contained in $N'$.  In particular task B.10(b), and the
corresponding rational-cusp implication in B.10(c) for
$N/N''$ and the marked quotient, are false.  This example has
only one ramified rational prime away from three; it makes
no assertion about equality of colours at distinct rational
primes or about infinitely many ramified primes.
\end{statement}
\begin{proof}
We first specify the based action, because using merely an
outer action would not suffice for this construction.
The cubic power cover of the tripod identifies the actual
geometric group with the characteristic open subgroup
\[
        N=\ker\bigl(F_2^{(3)}\longrightarrow(\mathbb Z/3)^2\bigr).
                                                               \tag{B.30.1}
\]
Its rank is $1+9(2-1)=10$.
At $O$ the lift of the standard tripod tangent is given by
$X/Z=-s$, $t=s^3$, $Y/Z=-(1-s^3)^{1/3}$.
Coefficientwise Galois action on the coherent Puiseux fibres
makes (B.30.1) an identification for the genuine based
arithmetic action, as proved in B.27.
Over $K$ the action on the Frattini quotient of $F_2^{(3)}$
is trivial: it is the mod-three cyclotomic scalar.
The kernel of the action on that Frattini quotient is pro-three,
by the finite three-group basis theorem applied to successive
characteristic finite quotients.
Its compact arithmetic subgroup restricts continuously to $N$.
The pair and the chosen unit tangent have good reduction
outside three, so inertia there acts trivially on $N$ by
prime-to-residue-characteristic specialization
\citep[XIII, Corollary 2.12]{BSGA1}; this is the based
specialization argument of B.27.
Consequently the actual action $\rho_N$ factors through
$\Gamma_S$ of (B.29.3), and its restriction to $P_S$ has
pro-three image.  This is a nonlinear statement about the
actual action on the entire free pro-three group, not only
about its abelianization.

The non-rational cusp $P$ is $K$-rational, and in the chosen
completion it is $\mathbb Q_7$-rational.  Choose its unit
tangent using the local parameter $X/Z$.
In the geometric pro-three fundamental exact sequence over
$\mathbb Q_7$, choose a geometric path from the fixed tangent
at $O$ to this tangent at $P$.  Its tangential section,
twisted inside its peripheral group by $a\kappa_3(7)$,
defines an actual continuous cocycle
\[
                z_{1,a}:G_{\mathbb Q_7}\longrightarrow N
\]
for the restriction of the original action $\rho_N$.
All paths are kept fixed as $a$ varies.  The local action has
pro-three image, so the graph homomorphism
\[
 g\longmapsto(z_{1,a}(g),\rho_N(g))
       \quad\text{in }N\rtimes\rho_N(G_{\mathbb Q_7})
                                                               \tag{B.30.2}
\]
has a pro-three target.  It factors through the actual maximal
pro-three quotient $L_1$.  We retain the same notation for its
cocycle there.  Under the literal local embedding in B.29,
$h\mapsto(z_{1,a}(h),h)$ is a continuous homomorphism
$L_1\to N\rtimes P_S$ whose second projection is the inclusion.

Choose the second local cocycle, on the conjugate factor,
by the exact formula
\[
                 z_{2,a}(chc)=\rho_N(c)z_{1,a}(h)
                         \qquad(h\in L_1).                    \tag{B.30.3}
\]
It is the actual local cusp section at $\bar P=cP$,
with the conjugate path and tangent.  The cocycle identity
follows from
$\rho_N(chc)=\rho_N(c)\rho_N(h)\rho_N(c)$.
We do not identify a positively normalized tame generator at
the second prime with $ct_1c$; complex conjugation acts on
the Tate twist and that normalization is unnecessary.

Both local graph sections now have target $N\rtimes P_S$,
a pro-three group.  By (B.29.1) and the universal property of
the free pro-three product they extend uniquely to
\[
            \sigma_a:P_S\longrightarrow N\rtimes P_S.
                                                               \tag{B.30.4}
\]
Its arithmetic projection is the identity, because it is
the identity on both generating local factors.  Write
$\sigma_a(h)=(z_a(h),h)$.  Formula (B.30.3) and uniqueness
give $z_a(chc)=\rho_N(c)z_a(h)$ on all of $P_S$.
Therefore setting $z_a(c)=1$ extends this to a continuous
cocycle on $\Gamma_S=P_S\rtimes\langle c\rangle$.
Let $r:G_{\mathbb Q}\to\Gamma_S$ be the actual quotient map.
The required section is
\[
                  s_a(g)=(z_a(r(g)),g).                       \tag{B.30.5}
\]
The second coordinate here is $g$, not $r(g)$.
Thus (B.30.5) is a genuine section of the original arithmetic
extension.  It is the action and the geometric cocycle which
factor through $\Gamma_S$, not the section's identity
arithmetic projection.  At every prime outside $\{3,7\}$
both the action and this cocycle are trivial on inertia.
Its localization at seven is the chosen local packet.

We finish with a conjugacy-invariant obstruction to global
cuspidality, already in $V=N^{\rm ab}$.
Let $D$ be the nine geometric coordinate cusps.  The usual
homology sequence of the punctured smooth cubic gives an
injection
\[
 W=\mathbb Z_3(1)[D]\big/
          \mathbb Z_3(1)\sum_{d\in D}e_d
           \lhook\joinrel\longrightarrow V.                  \tag{B.30.6}
\]
The generator of inertia about a cusp $d$ has image $e_d$.
At seven the unit tangential section at $P$ is unramified;
this follows from the smooth integral pair and its integral
unit tangent.  Choose $t_1$ so that $\kappa_3(7)(t_1)=1$.
Then the restriction of the abelian cocycle of $s_a$ to
this inertia generator is exactly
\[
                               a e_P.                        \tag{B.30.7}
\]
The geometric path only conjugates peripheral generators,
which has no effect in $V$.
Since inertia at seven acts trivially on $V$, a geometric
coboundary has zero restriction there.  Hence (B.30.7) is
invariant under global or local geometric conjugacy.
It also distinguishes all parameters $a$, since $e_P$
has infinite order.

The rational coordinate cusps are precisely
$[0:-1:1],[-1:0:1],[-1:1:0]$.
For a global rational cusp packet centred at such a cusp
$C$, its inertia value at seven belongs to
$\mathbb Z_3e_C$; changing a tangent merely changes that
scalar.  If $P\ne C$, the two classes $e_P,e_C$ are linearly
independent in (B.30.6): compare a third coordinate in
any alleged relation with the diagonal vector, then compare
the $P$ and $C$ coordinates.  Thus for $a\ne0$ equation
$ae_P=be_C$ is impossible for every $b\in\mathbb Z_3$.
This proves noncuspidality.

If $J\subseteq N'$, the quotient $N/J$ retains the entire
abelianization $V$.  The preceding obstruction therefore
survives its projection, as do ramification and distinctness
of parameters.  The intrinsic metabelian kernel $N''$ is
contained in $N'$.  In the exact marked notation of
LM7 E.71, $I_Q$ is the closed normal closure of a peripheral
generator $d_Q$ at the rational marked cusp $Q$, and
$J=\overline{\gamma_3(I_Q)}$.
This is normal in $N$ and Galois-stable: $I_Q$ is normal
and stable, and its closed lower central subgroups are
characteristic in it.  It is contained in $N'$, since its
defining triple commutators are.  Thus the same argument
applies to that specified marked quotient.
This transfer uses no assertion that
all sections of a quotient lift.
\end{proof}

\paragraph{Exact scope of the counterexample.}
The program \path{c_seven_free_product.py} independently
checks the Eisenstein norm identities, all residue-symbol
entries of (B.29.5), all $81$ possible upper-right blocks
in (B.29.6), and primitive independence of the relevant
boundary vectors.  The passage from these finite checks
to the full pro-three group is the cohomological proof
of B.29, not a computational inference.
If the chosen cusp in (B.30.2) is rational instead and the
local path is the restriction of a representative of its
global packet with value one at $c$, this conjugate gluing
recovers that global packet twisted by $a\kappa_3(7)$.
Indeed that global packet then has exactly the prescribed
restrictions to both local factors and the chosen value
at $c$, so uniqueness in (B.30.4) applies.
Without this compatibility on paths, local conjugacy is
not equality of the prescribed local graph maps, and no
global-cuspidality conclusion is inferred from uniqueness.
In either case this construction has only one possible
ramified rational prime away from three, so it does not
test equality of colours at different rational primes.
Tasks B.10(a) and the full pro-three form of B.11 remain
open in this return.  Adding an auxiliary rational prime
to the present construction introduces its actual
Frobenius word; the four-generator free-product proof
does not make that extra word independent.


\begin{statement}{B.31}{The entire zero-abelian slice has an exact linear order and decay law}{Proved}
Keep the $C_0$ tangent of B.28 and put $M=\Lambda^2V$.
Let $w=(3f_1+U_1)\wedge(3f_2+U_2)$, the primitive
integral Tate vector of B.7 and B.24.
The zero-abelian fibre in $H^1(\mathbb Q,N_2)$ is
$H^1(\mathbb Q,M)$, and every element has a unique
expression
\[
 z=z_0+\left(\sum_{p\ne3}a_p\kappa_3(p)\right)w,
 \quad z_0\in H^1(G_{\mathbb Q,\{3,\infty\}},M),
 \quad a_p\longrightarrow0.                                  \tag{B.31.1}
\]
Define the actual integral extension and its class by
\[
 \begin{gathered}
 0\longrightarrow A\longrightarrow
 L=\gamma_2N/\gamma_4N\longrightarrow M\longrightarrow0,\\
 \eta(g)=\chi_3(g)^{-1}g\omega-\omega\in A(-1),
       \qquad \omega\longmapsto w .
 \end{gathered}                                                 \tag{B.31.2}
\]
For a prime $\mathfrak q\mid p\ne3$ of $K$, put
\[
 \begin{gathered}
 q=N\mathfrak q,\qquad D_{\mathfrak q}=q^{-1}\Phi_{A,\mathfrak q}-I,
 \qquad
 e_{\mathfrak q}=[\eta(F_{\mathfrak q})]
                   \in A/D_{\mathfrak q}A,\\
 \operatorname{ord}(e_{\mathfrak q})=3^{m_p}.
 \end{gathered}                                                 \tag{B.31.3}
\]
The order is independent of the chosen prime over $p$,
prolongation, and integral lift $\omega$.
Then (B.31.1) lifts to intrinsic class three if and only if
\[
 \begin{split}
 v_3(a_p)&\ge m_p\quad\text{for every }p\ne3,\\
 v_3(a_p)-m_p&\longrightarrow+\infty
                    \quad\text{on the primes with }m_p>0.
 \end{split}                                                    \tag{B.31.4}
\]
Thus this restricted image is a subgroup. Its local-order
function is unbounded, and the image is a proper dense
subgroup of the locus satisfying only the first line of
(B.31.4). Its nonempty fibres are exactly B.28.

The integers $m_p$ are orders of specified finite actual
extension classes; they have not all been evaluated here.
B.35 proves $m_p\ge1$ at every $p\ne3$, so no unscaled
$\kappa_3(p)w$ lifts. B.36 evaluates the first actual
local order: $m_2=3$. Thus $a\kappa_3(2)w$ lifts exactly
when $27\mid a$. The smallest remaining local numerical
instance is $m_5$, with $1\le m_5\le5$ from the bounds below.
B.37 proves nonvanishing of the rational arithmetic Ross
component; it does not by itself evaluate these remaining
integral local orders.
This precise reduction is not a claim that the arbitrary
ramified-input image required in prompt B.5 has been evaluated.
\end{statement}
\begin{proof}
A class with zero abelianization can be geometrically
conjugated into the central kernel $M$ of $N_2\to V$.
The possible additional ambiguity is $V^{G_{\mathbb Q}}$,
which is zero by weights. This gives the stated
identification with $H^1(\mathbb Q,M)$.
At a good prime the only rational Frobenius eigenline of
$M$ of eigenvalue $p$ is the compact determinant line.
All other eigenvalues have absolute values $p^{3/2}$ or
$p^2$. The vector $w$ is primitive: its coefficient at
$A_0\wedge A_2$ is one. Every integral tame inertia value
is consequently $a_pw$ for a unique $a_p\in\mathbb Z_3$.
Continuity implies $a_p\to0$.
The uniformly convergent Kummer sum in (B.31.1) has exactly
these residues, so subtracting it leaves an unramified class.
This proves existence and uniqueness of (B.31.1).

The group $L$ is abelian because
$[\gamma_2N,\gamma_2N]\subset\gamma_4N$.
An ordered integral Hall section identifies it as a free
$\mathbb Z_3$-module extension; no BCH division by three
is used. The connecting homomorphism on $H^1(M)$ is
linear. Its value on $z_0$ vanishes by the finite-base
$H^2(A)$ vanishing in B.7. A cochain lift of
$\beta w$, where $\beta=\sum a_p\kappa_3(p)$, is
$g\mapsto\beta(g)\omega$. Its differential is precisely
$\chi_3(g)\beta(h)\eta(g)$, so its obstruction is
$[\eta\cup\beta]$.

The matrix $D_{\mathfrak q}$ is integral and rationally
invertible by the weights of $A(-1)$.
The actual tame correction is
\[
             -a_pD_{\mathfrak q}^{-1}\eta(F_{\mathfrak q}).
                                                               \tag{B.31.5}
\]
Its class modulo $A$ is killed by exactly $3^{m_p}$.
If $m_p>0$, the minimum valuation of the coordinates of
$D_{\mathfrak q}^{-1}\eta(F_{\mathfrak q})$ is exactly
$-m_p$. If $m_p=0$, that vector is integral; multiplication
by $a_p\to0$ then already supplies convergence.
This proves precisely the two conditions in (B.31.4).
The continuous localization isomorphism (B.24.19), applied
to $A$, proves necessity and sufficiency globally, including
the requirement of a single continuous lift.

Changing $\omega$ adds a coboundary $D_{\mathfrak q}b$
to $\eta(F_{\mathfrak q})$, so it does not change
$e_{\mathfrak q}$ or its order. A different prolongation
conjugates the module and this cohomology class integrally.
Conjugation over $\mathbb Q$ identifies the two primes
above a split $p$ and their orders. These observations
also prove independence of the choices in (B.31.3).

Here is an explicit finite matrix bound which avoids
replacing the local order by an unproved numerical value.
For an integral rationally invertible matrix $B$, define
the integer $\epsilon(B)$ by
$\operatorname{exp}(\operatorname{coker}B)=3^{\epsilon(B)}$.
Using the compact two-by-two matrix $A_E$ of (B.24.3), set
\[
 b_{\mathfrak q}=\max\left\{
 \begin{array}{l}
 \epsilon(A_E-I),\quad\epsilon(A_E\otimes A_E-I),\\
 \epsilon(qA_E-I),\quad v_3(q^2-1)
 \end{array}\right\}.
                                                               \tag{B.31.6}
\]
Then
\[
        m_p\le\epsilon(D_{\mathfrak q})
                    \le b_{\mathfrak q}+3.                    \tag{B.31.7}
\]
Indeed the integral map with columns $(3I_2,U)^{\rm t}$
and the eight cusp vectors splits $V$ rationally as
$H\oplus W$, and its inverse has denominator at most three.
The corresponding integral degree-three Lie map has
inverse denominator at most $27$.
In the split lattice the four multidegrees of $A(-1)$
are respectively $H$, eight copies of $H\otimes H$,
sixty-four copies of $H(1)$, and $168$ copies of
$\mathbb Z_3(2)$. Their Frobenius-minus-one matrices
are exactly the four matrices in (B.31.6).
Conjugating back proves (B.31.7).
Consequently the fully explicit sufficient conditions
\[
 v_3(a_p)\ge b_{\mathfrak q}+3,
 \qquad v_3(a_p)-b_{\mathfrak q}\longrightarrow+\infty
                                                               \tag{B.31.8}
\]
provide a lift without evaluating $\eta$.
For inert $p$, put $h=v_3(p+1)$. Here $q=p^2$ and
$A_E=-pI$, so $b_{\mathfrak q}=h+1$ by LTE.
In particular $v_3(a_p)\ge h+4$, with an additional
amount tending to infinity, suffices. This is a uniform
improvement of the coarse determinant bound in B.7.

At $(2)$ the fresh program \texttt{b\_class3\_fibres.py}
constructs $\Phi_A-4$ and certifies its complete Smith
valuations
\[
              0^{132},\quad1^{70},\quad2^{24},\quad3^{76},\quad4^{28}.
                                                               \tag{B.31.9}
\]
Their sum $458$ agrees independently with the eigenvalue
calculation $2+32+256+168$. Division by four does not
change them, so the local coefficient exponent is $3^4$.
This proves the bound on $m_2$, without pretending that the
coefficient-group exponent is the order of the actual class.
B.36 subsequently evaluates that class and gives $m_2=3$.

For unboundedness use the integral equivariant detector
$\Psi:A(-1)\to\mathbb Z_3(1)$ of B.17, for which
$\Psi_*[\eta]=6\kappa_3(3)$.
Take the primes $p_j$ of B.15:
$v_3(p_j-1)=j$ and three is not a cube modulo $p_j$.
The corresponding Kummer residue is a unit. Its multiple
by six in $\mathbb Z_3/(p_j-1)$ has order $3^{j-1}$.
Therefore
\[
                         m_{p_j}\ge j-1.                      \tag{B.31.10}
\]
Let $\mathcal L$ be the restricted coefficient families
$a_p\to0$ satisfying the first line of (B.31.4).
Finite-support families in $\mathcal L$ lift and are
dense in $\mathcal L$ for the uniform three-adic topology.
Along the displayed distinct primes, the family
$a_{p_j}=3^{m_{p_j}}$, zero elsewhere, belongs to
$\mathcal L$ but has correction valuation zero at every
$j\ge2$. It does not lift continuously. Thus the image
is proper and dense in the locally liftable locus.
B.17 supplies a different phantom with independently
specified numerical coefficients; the present family
proves the structure of the entire linear slice.
\end{proof}

\begin{statement}{B.32}{The inherited tripod metabelian quotient kills the compact determinant extension}{Proved}
Let $F$ be the geometric pro-three tripod group and let
$N=\ker(F\to(\mathbb Z/3)^2)$ be its actual Fermat cubic
cover subgroup. The inherited quotient $F/F''$ kills
the compact determinant extension used to define $\eta$.
Consequently its Jacobi-sum action does not directly
supply the missing intrinsic class-three elliptic value.
\end{statement}
\begin{proof}
The actual inclusion has exact abelian sequence
\[
 0\longrightarrow F'/N'\longrightarrow N^{\rm ab}
           \longrightarrow3F^{\rm ab}\longrightarrow0.
                                                               \tag{B.32.1}
\]
The last map is induced by the two coordinate units.
Both characters vanish on the rational compact summand
$H\subset V\otimes\mathbb Q_3$: their target is Tate
of weight two, whereas $H$ has weight one.
Thus the integral transfer vectors $u,v$ spanning the
compact lattice used above have representatives
$\widetilde u,\widetilde v\in F'$.
Their commutator $\omega=[\widetilde u,\widetilde v]$
lies in $F''$ and maps to $w=u\wedge v$ in
$\gamma_2N/\gamma_3N$.
Inside the abelian group $L=\gamma_2N/\gamma_4N$,
replace this lift by $(\omega-c\omega)/2$ if the
normalization $c\omega=-\omega$ is required.
The replacement remains in the image of $F''$.
Since $F''$ is characteristic in $F$, both $g\omega$
and every power of $\omega$ remain in $F''$.
Consequently the image of the actual difference
$g\omega\,\omega^{-\chi_3(g)}$, and hence of the
extension cocycle $\eta$, is zero under
\[
 N/\gamma_4N\longrightarrow N/(F''\gamma_4N).
                                                               \tag{B.32.2}
\]
This argument concerns the actual quotient map and its
kernel. It does not rule out a reconstruction theorem
using additional information, but such a theorem cannot
be replaced by reading the vanished extension in the
inherited metabelian quotient. The analogous free group
coming from $E\setminus\{O\}$ is a different ambient
cover and must not be substituted into (B.32.1).
\end{proof}
